\documentclass[
  oneside,
  openany,
  titlepage,
  11pt,
  a4paper,
  italian,
  DIV=14,
  headinclude,
  footinclude=true,
  cleardoublepage=empty
]{scrbook}

% --- ClassicThesis + ArsClassica (Pantieri) ---
\usepackage[T1]{fontenc}
\usepackage[utf8]{inputenc}
\usepackage{babel}
\usepackage{lmodern}
\usepackage{mathpazo}
\usepackage{amsmath,amssymb,amsthm}
\usepackage{mathtools}
\usepackage{mathrsfs}
\usepackage[
  eulerchapternumbers,
  beramono,
  eulermath,
  pdfspacing,
  parts
]{classicthesis}
\usepackage{arsclassica}

% --- nota di colore vinaccia per il numero del capitolo ---
\definecolor{halfgray}{rgb}{0.55,0.13,0.16}
\definecolor{futurered}{RGB}{180,18,18}
\definecolor{cream}{RGB}{245,240,228}

% --- aumenta lo spazio dopo il titolo del capitolo ---
\usepackage{titlesec}
\titlespacing*{\chapter}{0pt}{2.5\baselineskip}{4\baselineskip}

% --- lingua e tolleranze ---
\selectlanguage{italian}
\emergencystretch=3em
\tolerance=2000
\hbadness=2000

% --- grafica ---
\usepackage{graphicx}
\usepackage{tikz}
\usetikzlibrary{arrows.meta,positioning,decorations.pathmorphing,calc}

% --- epigrafi (per attacchi capitolo, opzionale) ---
\usepackage{epigraph}
\setlength{\epigraphwidth}{0.55\textwidth}
\renewcommand{\epigraphflush}{flushright}
\renewcommand{\textflush}{flushright}
\renewcommand{\epigraphsize}{\small\itshape}

% --- regole di inferenza in stile naturale ---
\usepackage{mathpartir}

% --- codice Agda con listings (no shell-escape richiesto) ---
\usepackage{listings}
\usepackage{xcolor}
\definecolor{agdaKw}{rgb}{0.55,0.13,0.16}   % maroon, come halfgray
\definecolor{agdaCom}{rgb}{0.45,0.45,0.45}
\definecolor{agdaStr}{rgb}{0.20,0.40,0.20}

% --- mapping unicode per Agda (così pdflatex digerisce i simboli) ---
\lstdefinelanguage{Agda}{
  keywords={data, where, with, module, open, import, public, private,
            renaming, hiding, using, let, in, do, postulate, abstract,
            primitive, mutual, infix, infixl, infixr, syntax, pattern,
            instance, record, field, constructor, Set, Type, Prop,
            forall, λ, lambda},
  sensitive=true,
  comment=[l]{--},
  morecomment=[s]{\{-}{-\}},
  string=[b]",
  basicstyle=\ttfamily\small,
  keywordstyle=\color{agdaKw}\bfseries,
  commentstyle=\color{agdaCom}\itshape,
  stringstyle=\color{agdaStr},
  columns=fullflexible,
  keepspaces=true,
  showstringspaces=false,
  literate=
    {𝔹}{{$\mathbb{B}$}}1
    {ℕ}{{$\mathbb{N}$}}1
    {≡}{{$\equiv$}}1
    {∀}{{$\forall$}}1
    {→}{{$\rightarrow$}}1
    {⊢}{{$\vdash$}}1
    {⊥}{{$\bot$}}1
    {⊤}{{$\top$}}1
    {⟦}{{$\llbracket$}}1
    {⟧}{{$\rrbracket$}}1
    {λ}{{$\lambda$}}1
    {Π}{{$\Pi$}}1
    {Σ}{{$\Sigma$}}1
    {α}{{$\alpha$}}1
    {β}{{$\beta$}}1
    {δ}{{$\delta$}}1
    {ε}{{$\varepsilon$}}1
    {ρ}{{$\rho$}}1
    {τ}{{$\tau$}}1
    {φ}{{$\varphi$}}1
    {₁}{{$_1$}}1
    {₂}{{$_2$}}1
    {⟨}{{$\langle$}}1
    {⟩}{{$\rangle$}}1
    {≢}{{$\not\equiv$}}1
    {∎}{{$\blacksquare$}}1
    {∷}{{$::$}}1
    {↑}{{$\uparrow$}}1
    {↓}{{$\downarrow$}}1
    {ℝ}{{$\mathbb{R}$}}1
    {ℤ}{{$\mathbb{Z}$}}1
    {è}{{\`e}}1
    {à}{{\`a}}1
    {ò}{{\`o}}1
    {ì}{{\`\i }}1
    {ù}{{\`u}}1
    {é}{{\'e}}1
    {È}{{\`E}}1
    {À}{{\`A}}1
    {Ì}{{\`I}}1
    {Ò}{{\`O}}1
    {Ù}{{\`U}}1
    {É}{{\'E}}1,
}

\lstnewenvironment{agda}[1][]{\lstset{language=Agda,#1}}{}
\newcommand{\agdainl}[1]{\lstinline[language=Agda,basicstyle=\ttfamily\normalsize]{#1}}

% --- ambienti teorema-like in stile ClassicThesis (identici al sorgente avit-io) ---
\newtheoremstyle{ctplain}%
  {0.7\baselineskip}{0.5\baselineskip}%
  {\itshape}{}%
  {\bfseries\scshape}{.}%
  {.5em}{}%
\newtheoremstyle{ctdef}%
  {0.7\baselineskip}{0.5\baselineskip}%
  {\normalfont}{}%
  {\bfseries\scshape}{.}%
  {.5em}{}%
\newtheoremstyle{ctrem}%
  {0.7\baselineskip}{0.5\baselineskip}%
  {\normalfont}{}%
  {\itshape}{.}%
  {.5em}{}%

% --- ambiente per i ceri: stesso stile sobrio dello stile def, ma con head dedicato ---
\newtheoremstyle{ctcero}%
  {0.9\baselineskip}{0.7\baselineskip}%
  {\normalfont}{}%
  {\bfseries\scshape\color{halfgray}}{.}%
  {.5em}{}%

\theoremstyle{ctplain}
\newtheorem{regola}{Regola}[chapter]

\theoremstyle{ctdef}
\newtheorem{defn}[regola]{Definizione}

\theoremstyle{ctrem}
\newtheorem*{oss}{Osservazione}

\theoremstyle{ctcero}
\newtheorem{cero}{Cero}

% --- comandi locali ---
\newcommand{\meta}[1]{\textit{#1}}        % per il metalinguaggio in prosa
\newcommand{\lang}[1]{\textbf{#1}}        % per il linguaggio in prosa
\newcommand{\definitorio}{\ensuremath{\equiv}}
\newcommand{\identita}{\ensuremath{\equiv}}
\newcommand{\Type}{\mathsf{Type}}
\newcommand{\Set}{\mathsf{Set}}
\newcommand{\indB}{\mathsf{ind\text{-}\mathbb{B}}}
\newcommand{\indN}{\mathsf{ind\text{-}\mathbb{N}}}
\newcommand{\indS}{\mathsf{ind\text{-}\Sigma}}
\newcommand{\indV}{\mathsf{ind\text{-}Vec}}
\newcommand{\Jelim}{\mathsf{J}}
\newcommand{\refl}{\mathsf{refl}}
\newcommand{\zero}{\mathsf{zero}}
\newcommand{\suc}{\mathsf{suc}}
\newcommand{\tto}{\mathsf{tt}}
\newcommand{\ff}{\mathsf{ff}}
\newcommand{\nil}{\mathsf{nil}}
\newcommand{\cons}{\mathsf{cons}}
\newcommand{\andd}{\mathbin{\&\&}}

% --- hyperref ---
\hypersetup{%
  colorlinks=true,
  linkcolor=Maroon,
  citecolor=Maroon,
  urlcolor=Maroon,
  pdftitle={Prova o taci},
  pdfsubject={appunti di MLTT per Verified Functional Programming in Agda},
  pdfauthor={}
}

% --- TOC ---
\usepackage{titletoc}
\renewcommand{\cftpartfont}{\color{CTtitle}}
\renewcommand{\cftpartaftersnum}{}
\renewcommand{\cftpartaftersnumb}{}
\renewcommand{\cftchapfont}{}
\renewcommand{\cftchappresnum}{}
\renewcommand{\cftchapaftersnumb}{}

\begin{document}
\frenchspacing
\raggedbottom
\selectlanguage{italian}
\renewcommand{\partname}{Parte}
\renewcommand{\chaptername}{Capitolo}
\renewcommand{\contentsname}{Indice}

\pagenumbering{roman}
\pagestyle{plain}

% ----------------- frontespizio (stile Depero — Il Verificatore) -----------------
\begin{titlepage}
\thispagestyle{empty}
\begin{tikzpicture}[remember picture, overlay]

  %% ======= ZONE DI SFONDO =======

  \fill[black]
    (current page.north west) rectangle
    ($(current page.north east)+(0,-11.5cm)$);

  \fill[futurered]
    ($(current page.north west)+(0,-11.5cm)$) rectangle
    ($(current page.north east)+(0,-15cm)$);

  \fill[cream]
    ($(current page.north west)+(0,-15cm)$) rectangle
    ($(current page.north east)+(0,-24.5cm)$);

  \fill[black]
    ($(current page.north west)+(0,-24.5cm)$) rectangle
    (current page.south east);

  %% ======= TORRE SINISTRA =======

  \fill[futurered!55!black]
    ($(current page.north)+(-9.5cm,-1.5cm)$) rectangle
    ($(current page.north)+(-6.3cm,-11.5cm)$);
  %% finestre
  \foreach \yw in {2.8,4.8,6.8,8.8} {
    \fill[black!85]
      ($(current page.north)+(-9.1cm,-\yw cm)$) rectangle +(0.75cm,-0.9cm);
    \fill[black!85]
      ($(current page.north)+(-7.75cm,-\yw cm)$) rectangle +(0.75cm,-0.9cm);
  }
  %% cornicione torre sinistra
  \fill[cream!40!black]
    ($(current page.north)+(-9.7cm,-1.2cm)$) rectangle
    ($(current page.north)+(-6.1cm,-1.5cm)$);

  %% ======= TORRE DESTRA =======

  \fill[cream!45!black]
    ($(current page.north)+(6.3cm,-2cm)$) rectangle
    ($(current page.north)+(9.5cm,-11.5cm)$);
  %% finestre
  \foreach \yw in {3.3,5.3,7.3,9.3} {
    \fill[black!85]
      ($(current page.north)+(6.7cm,-\yw cm)$) rectangle +(0.75cm,-0.9cm);
    \fill[black!85]
      ($(current page.north)+(8.05cm,-\yw cm)$) rectangle +(0.75cm,-0.9cm);
  }
  %% cornicione torre destra
  \fill[cream!40!black]
    ($(current page.north)+(6.1cm,-1.7cm)$) rectangle
    ($(current page.north)+(9.7cm,-2cm)$);

  %% ======= IL VERIFICATORE =======

  %% -- Cappello: triangolo crema con Π --
  \fill[cream]
    ($(current page.north)+(-2.1cm,-1.3cm)$) --
    ($(current page.north)+(2.1cm,-1.3cm)$) --
    ($(current page.north)+(0cm,-3.7cm)$) -- cycle;
  \draw[black, line width=2.5pt]
    ($(current page.north)+(-2.1cm,-1.3cm)$) --
    ($(current page.north)+(2.1cm,-1.3cm)$) --
    ($(current page.north)+(0cm,-3.7cm)$) -- cycle;
  \node[black, font=\fontsize{20}{20}\selectfont\bfseries]
    at ($(current page.north)+(0,-3cm)$) {$\Pi$};

  %% -- Testa --
  \fill[cream]
    ($(current page.north)+(-1.05cm,-3.7cm)$) rectangle
    ($(current page.north)+(1.05cm,-5.45cm)$);
  \draw[black, line width=2.5pt]
    ($(current page.north)+(-1.05cm,-3.7cm)$) rectangle
    ($(current page.north)+(1.05cm,-5.45cm)$);
  %% occhi
  \fill[futurered]
    ($(current page.north)+(-0.42cm,-4.55cm)$) circle (0.22cm);
  \fill[futurered]
    ($(current page.north)+(0.42cm,-4.55cm)$) circle (0.22cm);
  %% bocca
  \draw[black, line width=2pt]
    ($(current page.north)+(-0.38cm,-5.1cm)$) --
    ($(current page.north)+(0.38cm,-5.1cm)$);

  %% -- Collo --
  \fill[black!55]
    ($(current page.north)+(-0.32cm,-5.45cm)$) rectangle
    ($(current page.north)+(0.32cm,-5.95cm)$);
  \draw[black, line width=1.5pt]
    ($(current page.north)+(-0.32cm,-5.45cm)$) rectangle
    ($(current page.north)+(0.32cm,-5.95cm)$);

  %% -- Torso sinistro (rosso futurista) --
  \fill[futurered]
    ($(current page.north)+(-2.3cm,-5.95cm)$) rectangle
    ($(current page.north)+(0cm,-9.1cm)$);

  %% -- Torso destro (crema) --
  \fill[cream]
    ($(current page.north)+(0cm,-5.95cm)$) rectangle
    ($(current page.north)+(2.3cm,-9.1cm)$);

  %% -- Bordo e divisore torso --
  \draw[black, line width=2.5pt]
    ($(current page.north)+(-2.3cm,-5.95cm)$) rectangle
    ($(current page.north)+(2.3cm,-9.1cm)$);
  \draw[black, line width=2pt]
    ($(current page.north)+(0,-5.95cm)$) --
    ($(current page.north)+(0,-9.1cm)$);

  %% -- Simboli sul torso --
  \node[white, font=\fontsize{22}{22}\selectfont]
    at ($(current page.north)+(-1.15cm,-7.5cm)$) {$\Sigma$};
  \node[futurered!60!black, font=\fontsize{18}{18}\selectfont\bfseries]
    at ($(current page.north)+(1.15cm,-7.5cm)$) {$\equiv$};

  %% -- Spallaccio sinistro --
  \fill[futurered!70!black]
    ($(current page.north)+(-2.3cm,-5.95cm)$) rectangle
    ($(current page.north)+(-1.3cm,-6.55cm)$);

  %% -- Spallaccio destro --
  \fill[cream!60!black]
    ($(current page.north)+(1.3cm,-5.95cm)$) rectangle
    ($(current page.north)+(2.3cm,-6.55cm)$);

  %% -- Braccio sinistro --
  \fill[futurered!75!black]
    ($(current page.north)+(-2.3cm,-6.55cm)$) rectangle
    ($(current page.north)+(-4.7cm,-7.35cm)$);
  \draw[black, line width=2pt]
    ($(current page.north)+(-2.3cm,-6.55cm)$) rectangle
    ($(current page.north)+(-4.7cm,-7.35cm)$);

  %% -- Mano sinistra: cerchio con λ --
  \fill[cream]
    ($(current page.north)+(-5.45cm,-6.95cm)$) circle (0.75cm);
  \draw[black, line width=2pt]
    ($(current page.north)+(-5.45cm,-6.95cm)$) circle (0.75cm);
  \node[black, font=\fontsize{17}{17}\selectfont]
    at ($(current page.north)+(-5.45cm,-6.95cm)$) {$\lambda$};

  %% -- Braccio destro --
  \fill[cream!65!black]
    ($(current page.north)+(2.3cm,-6.55cm)$) rectangle
    ($(current page.north)+(4.7cm,-7.35cm)$);
  \draw[black, line width=2pt]
    ($(current page.north)+(2.3cm,-6.55cm)$) rectangle
    ($(current page.north)+(4.7cm,-7.35cm)$);

  %% -- Mano destra: cerchio con ⊢ --
  \fill[futurered]
    ($(current page.north)+(5.45cm,-6.95cm)$) circle (0.75cm);
  \draw[black, line width=2pt]
    ($(current page.north)+(5.45cm,-6.95cm)$) circle (0.75cm);
  \node[white, font=\fontsize{17}{17}\selectfont\bfseries]
    at ($(current page.north)+(5.45cm,-6.95cm)$) {$\vdash$};

  %% -- Cintura --
  \fill[black]
    ($(current page.north)+(-2.3cm,-9.1cm)$) rectangle
    ($(current page.north)+(2.3cm,-9.6cm)$);

  %% -- Gamba sinistra --
  \fill[futurered]
    ($(current page.north)+(-1.75cm,-9.6cm)$) rectangle
    ($(current page.north)+(-0.55cm,-11.25cm)$);
  \draw[black, line width=2pt]
    ($(current page.north)+(-1.75cm,-9.6cm)$) rectangle
    ($(current page.north)+(-0.55cm,-11.25cm)$);

  %% -- Gamba destra --
  \fill[cream!25!black]
    ($(current page.north)+(0.55cm,-9.6cm)$) rectangle
    ($(current page.north)+(1.75cm,-11.25cm)$);
  \draw[black, line width=2pt]
    ($(current page.north)+(0.55cm,-9.6cm)$) rectangle
    ($(current page.north)+(1.75cm,-11.25cm)$);

  %% -- Piede sinistro --
  \fill[black]
    ($(current page.north)+(-2.25cm,-11.25cm)$) rectangle
    ($(current page.north)+(-0.3cm,-11.52cm)$);

  %% -- Piede destro --
  \fill[black]
    ($(current page.north)+(0.3cm,-11.25cm)$) rectangle
    ($(current page.north)+(2.4cm,-11.52cm)$);

  %% ======= TITOLO (fascia rossa) =======

  \node[white, font=\fontsize{33}{33}\selectfont\bfseries\scshape]
    at ($(current page.north)+(0,-13.25cm)$)
    {\spacedallcaps{Prova o taci}};

  %% ======= SEZIONE CREMA =======

  %% regola di inferenza
  \node[black!68, font=\fontsize{13}{15}\selectfont]
    at ($(current page.north)+(0,-18.8cm)$) {$a : A$};
  \draw[black!55, line width=2.5pt]
    ($(current page.north)+(-3.6cm,-19.7cm)$) --
    ($(current page.north)+(3.6cm,-19.7cm)$);
  \node[black!68, font=\fontsize{13}{15}\selectfont]
    at ($(current page.north)+(0,-20.7cm)$)
    {$\mathsf{refl} : a \equiv a$};

  %% etichetta F·I·E·C
  \node[futurered, font=\fontsize{11}{13}\selectfont\bfseries\scshape]
    at ($(current page.north)+(0,-22.5cm)$)
    {F \quad I \quad E \quad C};

  %% ======= SOTTOTITOLO (fondo nero) =======

  \node[white!60!black, font=\fontsize{9}{11}\selectfont\itshape]
    at ($(current page.south)+(0,3.8cm)$)
    {compendio a \textit{Verified Functional Programming in Agda}};
  \node[futurered!65!white, font=\fontsize{8}{10}\selectfont]
    at ($(current page.south)+(0,2.2cm)$)
    {Aaron Stump \textbullet\ Iowa Agda Library};

\end{tikzpicture}
\end{titlepage}
\clearpage

% ----------------- nota al lettore -----------------
\chapter*{Nota al lettore}
\addcontentsline{toc}{chapter}{Nota al lettore}

\begin{center}
{\large\bfseries\scshape\color{futurered}
Nessuna magia.\\[2pt]
Nessun trucco.\\[2pt]
Il tipo esiste o non esiste.\\[2pt]
La prova regge o crolla.\\[2pt]
Chi non dimostra, tace.}
\end{center}

\bigskip\bigskip

Quattro cose, prima di cominciare.

\paragraph{Origine.} Queste note nascono per accompagnare la lettura di \emph{Verified Functional Programming in Agda} di Aaron Stump, leggendo gli esercizi come applicazioni della Martin-L\"of Type Theory pura --- cio\`e nello schema delle quattro regole F-I-E-C (Formazione, Introduzione, Eliminazione, Computazione) che sotto Agda c'\`e sempre, anche quando la sintassi la nasconde dietro lo zucchero.

\paragraph{Tesi.} Tutto il gioco di MLTT \`e in quattro regole, applicate uniformemente a ogni tipo. Capire ciascun tipo significa scriverne le quattro regole e vedere come gli abitanti del linguaggio (termini, prove) vengono dalle regole di introduzione, e come gli usi (eliminazioni, riduzioni) vengono da quelle di eliminazione e computazione. Non c'\`e altro.

\paragraph{Assunzioni.} Sono dichiarate ed esplicite. Le ho chiamate \emph{ceri} --- come i ceri votivi che si accendono perch\'e qualcosa regga senza averne la dimostrazione interna --- e ne incontreremo quattro nel corpo del testo. Ognuno \`e un punto in cui il sistema poggia su qualcosa che non si pu\`o derivare dall'interno: vivono nel metalinguaggio o sono scelte fondative. Confondere un cero con un teorema, o uno zucchero sintattico con una regola primitiva, \`e la fonte principale di confusione in MLTT. Renderli espliciti \`e met\`a del lavoro.

\paragraph{Forma.} Dove la prosa basta, c'\`e prosa. Dove la prosa rende sfuocato qualcosa che chiede precisione, entrano le regole di inferenza in forma classica (premesse, barra, conclusione), o un frammento Agda. Nessun salto: ogni sezione usa solo cose introdotte in sezioni precedenti, e se ci si perde su un punto basta tornare indietro fino a trovarne il fondamento.

\bigskip

\noindent\textbf{Patto di lettura.} Niente salti in avanti. La Parte~I fissa il vocabolario (giudizi, metalinguaggio, schema F-I-E-C). La Parte~II applica F-I-E-C a sei tipi, in ordine di complessit\`a crescente: $\mathbb{B}$, $\mathbb{N}$, $\Pi$, $\Sigma$, il tipo identit\`a $\_\identita\_$, $\mathrm{Vec}$. La Parte~III mostra come questi sei tipi, combinati, governano il primo esercizio non banale del libro di Stump (commutativit\`a di \texttt{\&\&}). I ceri e le decisioni di design sono raccolti in fondo (\S\ref{cap:ceri}).

\bigskip

\noindent\textbf{Strumenti.} Stile tipografico Pantieri (\textsc{ClassicThesis} + \textsc{ArsClassica}); ambienti teorema \textsc{regola} (numerate per capitolo), \textsc{definizione} (sullo stesso contatore), \textsc{osservazione} (senza numero), \textsc{cero} (numerato a parte, fuori dai capitoli, per tenere le assunzioni globali leggibili come un indice). Il codice Agda \`e reso con \texttt{listings} e mappa Unicode --- quindi compila con \texttt{pdflatex} senza shell-escape.

\clearpage

% ----------------- indice -----------------
{\setlength{\parskip}{0pt}\tableofcontents}
\clearpage

\pagenumbering{arabic}
\pagestyle{scrheadings}

% ----------------- corpo del libro -----------------
% =====================================================================
% PARTE I — Le fondamenta
% =====================================================================
\part{Le fondamenta}

\noindent\itshape Cosa assumi prima di scrivere la prima regola.
Cinque capitoli che fissano il vocabolario: il rapporto tra formalismo e materia
(\S\ref{cap:partenza}); i quattro giudizi del metalinguaggio (\S\ref{cap:giudizi});
la grammatica di contesti e dell'operazione $\vdash$ (\S\ref{cap:contesti});
il test operativo per riconoscere linguaggio e metalinguaggio
(\S\ref{cap:test}); lo schema F-I-E-C che governa ogni tipo
(\S\ref{cap:fiec}). Da qui in poi, nulla si introduce senza una di queste
cinque lavagne sotto.\upshape

\bigskip

% ---------------------------------------------------------------------
\chapter{Il punto di partenza}\label{cap:partenza}

{\spacedlowsmallcaps{Il formalismo fissa qualcosa che era gi\`a l\`i}}, non crea, cattura. Le regole non rendono vere le cose --- le cose sono vere indipendentemente dalle regole. Le regole sono uno strumento per essere espliciti su \emph{cosa si assume} e \emph{cosa si deriva}.

Questa distinzione \`e l'asse di tutto il libro. Ogni volta che il terreno smette di essere derivabile e diventa assunto, il testo lo segnala con un \emph{cero}: una proposizione numerata, scritta in stile sobrio, che enuncia l'assunzione e dice perch\'e regge anche senza dimostrazione interna. I ceri sono raccolti in fondo (\S\ref{cap:ceri}) e si possono leggere in sequenza per avere il riassunto completo delle assunzioni di MLTT.

\bigskip

\noindent La metafora del cero votivo non \`e ornamento. Un cero si accende perch\'e qualcosa --- non importa cosa --- regga: l'atto di accenderlo non costituisce prova, e tuttavia non \`e capriccio, perch\'e ci\`o che si \`e affidato al cero ha un peso reale. Nel formalismo accade lo stesso. Certi giudizi --- l'esistenza del metalinguaggio, l'esistenza dei contesti, lo statuto delle regole di computazione, la decisione di lasciare entrare termini dentro i tipi --- non si possono dimostrare dall'interno del sistema. Si accendono come ceri: si dichiarano, si tengono in vista, e si va avanti.

\bigskip

\noindent\textbf{Cosa non sono i ceri.} Non sono assiomi nel senso di Hilbert (proposizioni del linguaggio scelte come punto di partenza per derivare altre proposizioni del linguaggio): vivono nel metalinguaggio, o sono scelte di design. Non sono congetture (non si tratta di affermazioni che potrebbero un giorno rivelarsi false). Non sono postulati arbitrari: sono i punti in cui la materia del sistema chiede di essere fissata, e ogni fissaggio \`e una scelta motivata ma non costretta.

\bigskip

\noindent\textbf{Quanti sono.} Quattro nel corpo del libro, pi\`u un quinto al limite (la coerenza del sistema, indimostrabile dall'interno per Gödel). Sono pochi perch\'e MLTT \`e parsimoniosa. Una volta accesi, le quattro regole F-I-E-C (\S\ref{cap:fiec}) generano l'intero edificio.

% ---------------------------------------------------------------------
\chapter{I quattro giudizi del metalinguaggio}\label{cap:giudizi}

Tutto vive qui, fuori dal linguaggio oggetto. I quattro giudizi sono:

\begin{equation*}
\begin{array}{ll}
A : \Type & \text{\meta{$A$ \`e un tipo}} \\
a : A & \text{\meta{$a$ abita $A$}} \\
A \definitorio B : \Type & \text{\meta{$A$ e $B$ sono lo stesso tipo (uguaglianza definitoria)}} \\
a \definitorio b : A & \text{\meta{$a$ e $b$ sono lo stesso termine in $A$ (uguaglianza definitoria)}}
\end{array}
\end{equation*}

Questi \emph{non sono} tipi o termini: sono \textbf{affermazioni sul sistema}. Vivono nel metalinguaggio. Prima ancora di scrivere una regola, stai gi\`a assumendo che questo livello esista e abbia senso.

\begin{cero}[Esiste un metalinguaggio con i quattro giudizi]
Non lo si dimostra dall'interno del sistema: \`e la lavagna su cui si scrivono le regole. Prima assunzione, quella che rende possibili tutte le altre.
\end{cero}

\begin{cero}[Esistono i contesti e l'operazione $\vdash$]\label{cero:contesti}
``Assumere $x : A$'' \`e un'operazione legittima: i contesti si possono estendere, si pu\`o ragionare \emph{sotto} un'ipotesi. Senza, niente regole con premesse ipotetiche, niente tipi dipendenti.
\end{cero}

\begin{oss}[Sulle due uguaglianze]
L'uguaglianza $\definitorio$ di questa sezione \`e \textbf{definitoria}: un'asserzione sul sistema, non un tipo, non si dimostra. \`E quella che dice ``queste due cose riducono alla stessa cosa per le regole di computazione''. La incontreremo costantemente da qui a tutto il libro.

In \S\ref{cap:identita} introdurremo un \emph{altro} $\identita$, completamente diverso, che vive \textbf{nel linguaggio}: \`e un tipo, lo si abita con prove, lo si manipola come qualsiasi altro termine. Nei testi classici si scrivono uguali ($\equiv$ in entrambi i casi), e questa \`e una sciagura notazionale storica; in questo libro li distinguiamo tipograficamente apposta: $a \definitorio b : A$ per il primo, $a \identita b$ per il secondo, oppure $\mathsf{Id}\, A\, a\, b$ se serve massima chiarezza.
\end{oss}

% ---------------------------------------------------------------------
\chapter{Contesti e l'operazione $\vdash$}\label{cap:contesti}

\noindent {\spacedlowsmallcaps{Tutto quello che scriveremo da qui in poi vale \emph{sotto ipotesi}}}. Non c'\`e regola in MLTT --- nemmeno la formazione di $\mathbb{B}$ --- che non si svolga dentro un contesto. Era il Cero~\ref{cero:contesti}; resta da specificarne la grammatica.

\bigskip

\noindent\textbf{Il caso minimo dove serve.} Vogliamo poter scrivere la regola di formazione di $\Pi$ (\S\ref{cap:pi}):
\[
\inferrule{A : \Type \\ x : A \vdash B\,x : \Type}{(x : A) \to B\,x : \Type}
\]
La seconda premessa dice ``in un contesto in cui $x : A$ \`e assunto, $B\,x$ \`e un tipo''. Senza una nozione formale di contesto e di operazione $\vdash$, quella riga \`e gergo. Con la grammatica esplicita, \`e una regola.

\section*{Grammatica dei contesti}

I contesti sono sequenze finite di ipotesi della forma $x : A$. Li indichiamo con $\Gamma, \Delta$, eventualmente con apici. Introduciamo un nuovo giudizio del metalinguaggio:
\[
\Gamma \;\mathsf{ctx} \qquad \text{\meta{$\Gamma$ \`e un contesto ben formato}}
\]

\noindent Le regole di formazione del giudizio $\mathsf{ctx}$ sono due, e generano tutti e soli i contesti legittimi.

\paragraph{Contesto vuoto.}
\[
\inferrule{ }{\cdot \;\mathsf{ctx}}
\]
Il contesto vuoto, sempre ben formato.

\paragraph{Estensione.}
\[
\inferrule{\Gamma \;\mathsf{ctx} \\ \Gamma \vdash A : \Type \\ x \notin \mathrm{dom}(\Gamma)}{\Gamma, x : A \;\mathsf{ctx}}
\]
Si pu\`o estendere $\Gamma$ con un'ipotesi $x : A$ \emph{a patto che} $A$ sia un tipo derivabile in $\Gamma$ (terza premessa) e $x$ sia un nome nuovo, non gi\`a presente (quarta premessa, condizione di freschezza).

\bigskip

\noindent\textbf{Osservazioni meccaniche.}
\begin{itemize}
\item Il contesto cresce \emph{a destra}: $\Gamma, x : A$ si legge ``$\Gamma$ esteso con $x : A$''. Le ipotesi pi\`u vecchie stanno a sinistra.
\item Le ipotesi pi\`u recenti possono dipendere dalle precedenti --- il tipo $B\,x$ in $\Gamma, x : A, y : B\,x$ usa la $x$ del passo prima. Senza questa possibilit\`a non si avrebbero tipi dipendenti, ed \`e proprio questa l'apertura formale del Cero~\ref{cero:dipendenza} che vedremo in \S\ref{cap:pi}.
\item Il giudizio $\mathsf{ctx}$ \`e del \textbf{metalinguaggio}, come gli altri quattro; vive sopra il linguaggio.
\end{itemize}

\section*{I quattro giudizi, riscritti sotto contesto}

I quattro giudizi di \S\ref{cap:giudizi} non compaiono mai \emph{nudi} nelle regole reali: portano sempre un contesto a sinistra di $\vdash$.

\begin{equation*}
\begin{array}{ll}
\Gamma \vdash A : \Type & \text{\meta{nel contesto $\Gamma$, $A$ \`e un tipo}} \\
\Gamma \vdash a : A & \text{\meta{nel contesto $\Gamma$, $a$ abita $A$}} \\
\Gamma \vdash A \definitorio B : \Type & \text{\meta{$A$ e $B$ sono lo stesso tipo in $\Gamma$}} \\
\Gamma \vdash a \definitorio b : A & \text{\meta{$a$ e $b$ sono lo stesso termine di $A$ in $\Gamma$}}
\end{array}
\end{equation*}

\noindent Il contesto non \`e ``parte'' del giudizio: \`e \emph{l'ambiente di ipotesi sotto cui il giudizio vale}. Contesti diversi producono giudizi diversi della stessa forma: $\cdot \vdash a : A$ (sotto vuoto) e $x : B \vdash a : A$ sono due asserzioni distinte, anche con $a$ e $A$ identici.

\section*{La regola della variabile}

Senza una regola che permetta di \emph{usare} le ipotesi, il contesto sarebbe decorativo. La regola che le mette in gioco \`e la pi\`u semplice possibile:

\[
\inferrule{\Gamma, x : A, \Delta \;\mathsf{ctx}}{\Gamma, x : A, \Delta \vdash x : A}
\]

\noindent Letta: ``se nel contesto compare $x : A$, allora posso derivare $x : A$''. \`E ovvia; \`e quella che giustifica formalmente l'idea ``ho $x$ a disposizione perch\'e l'ho assunto''. \`E anche l'unica regola che ``guarda dentro'' il contesto --- tutte le altre lo lasciano com'\`e o lo estendono in modo prevedibile.

\section*{Le altre regole strutturali}

Tre regole che governano \emph{come} il contesto pu\`o cambiare lungo una derivazione senza alterarne il significato. Le scriviamo nella loro forma minima; lo scope della trattazione le tiene a questo livello: \emph{The Little Typer} e \emph{Verified Functional Programming in Agda} non chiedono di pi\`u. Nessuna di esse \`e un cero: sono \emph{lemmi ammissibili}, conseguenze del modo in cui le regole specifiche dei tipi (Parte~II) saranno scritte. Si dichiarano qui perch\'e da Parte~II in poi le useremo tacitamente.

\paragraph{Indebolimento (\emph{weakening}).}
\[
\inferrule{\Gamma, \Delta \vdash \mathcal{J} \\ \Gamma \vdash A : \Type \\ x \notin \mathrm{dom}(\Gamma, \Delta)}{\Gamma, x : A, \Delta \vdash \mathcal{J}}
\]
Aggiungere un'ipotesi non rovina ci\`o che era gi\`a derivato: se sapevi $\mathcal{J}$, lo sai ancora con un'ipotesi in pi\`u. Qui $\mathcal{J}$ sta per uno qualsiasi dei quattro giudizi --- usiamo la calligrafica $\mathcal{J}$ per non confonderla con l'eliminatore $\Jelim$ del tipo identit\`a (\S\ref{cap:identita}).

\paragraph{Sostituzione.}
\[
\inferrule{\Gamma \vdash a : A \\ \Gamma, x : A, \Delta \vdash \mathcal{J}}{\Gamma, \Delta[a/x] \vdash \mathcal{J}[a/x]}
\]
Se in $\Gamma$ esibisci un abitante $a$ di $A$, ogni derivazione fatta \emph{sotto l'ipotesi $x : A$} si pu\`o specializzare a quel concreto $a$, sostituendolo dappertutto --- nelle ipotesi successive che dipendevano da $x$ e nel giudizio finale. \`E la regola che chiude l'ipotesi.

\paragraph{Conversione di tipo.}
\[
\inferrule{\Gamma \vdash a : A \\ \Gamma \vdash A \definitorio B : \Type}{\Gamma \vdash a : B}
\]
Se $A$ e $B$ sono definitoriamente uguali, abitare $A$ e abitare $B$ \`e la stessa cosa. \`E la regola che fa lavorare $\definitorio$ \emph{dentro} il sistema dei tipi: senza, l'uguaglianza definitoria si fermerebbe a essere un'asserzione esterna e non avrebbe effetto operativo sul controllo dei tipi.

\section*{Convenzione tipografica per il resto del libro}

Da \S\ref{cap:fiec} in poi, le regole di inferenza si scrivono \emph{senza contesto in vista} quando tutte le premesse e la conclusione condividono lo stesso. Si sottintende che ce ne sia uno, lo stesso ovunque nella regola. Esempio: la regola di introduzione di $\mathbb{B}$ (\S\ref{cap:bool})
\[
\inferrule{ }{\tto : \mathbb{B}}
\]
sta in realt\`a per
\[
\inferrule{\Gamma \;\mathsf{ctx}}{\Gamma \vdash \tto : \mathbb{B}}
\]
per qualunque $\Gamma$ ben formato.

\bigskip

\noindent Quando invece una premessa \emph{estende} il contesto rispetto alla conclusione --- caso paradigmatico, la $\Pi$-formazione di \S\ref{cap:pi}:
\[
\inferrule{A : \Type \\ x : A \vdash B\,x : \Type}{(x : A) \to B\,x : \Type}
\]
--- il $\vdash$ sopravvive nella scrittura, perch\'e \`e l\`i che il contesto sta facendo un lavoro non banale. La regola, sotto contesto generico $\Gamma$, sarebbe per esteso:
\[
\inferrule{\Gamma \vdash A : \Type \\ \Gamma, x : A \vdash B\,x : \Type}{\Gamma \vdash (x : A) \to B\,x : \Type}
\]
La compressione ``ometti $\Gamma$ dove \`e ovvio, lascialo dove non lo \`e'' \`e un fatto di leggibilit\`a; sotto, la regola \`e la seconda.

\section*{Test diretto}

\begin{oss}[$\vdash$ e ``,'' sono entrambi del metalinguaggio]
Il ``,'' che separa le ipotesi in $\Gamma, x : A$ \`e grammatica dei contesti, non del linguaggio: non puoi passare $\Gamma, x : A$ come argomento di una funzione, perch\'e non \`e un oggetto del linguaggio. \`E un contesto, e i contesti vivono nel metalinguaggio. Stessa cosa per $\vdash$ --- gi\`a sotto test in \S\ref{cap:test}, ma ora con grammatica precisa alle spalle.
\end{oss}

\bigskip

\noindent Con la grammatica di $\vdash$ in tasca, ogni regola del seguito ha un significato preciso e non interpretativo. Ogni volta che leggi una premessa con un nuovo $x : A$, sai cosa si sta facendo: estensione di contesto via la regola di estensione, e da quel momento $x$ \`e disponibile come variabile via la regola della variabile.

% ---------------------------------------------------------------------
\chapter{Linguaggio e metalinguaggio --- il test}\label{cap:test}

Sezione operativa: ogni volta che sui livelli ci si perde, si torna qui.

\section*{Il test in una riga}

\noindent\itshape Questo oggetto lo posso passare come argomento a una funzione, o ci posso costruire sopra un altro termine?\upshape

\medskip

\noindent \textbf{S\`i} $\Rightarrow$ vive nel \textbf{linguaggio}. \`E un termine o un tipo: oggetto manipolabile.

\noindent \textbf{No} --- \`e qualcosa che \emph{autorizza} a scrivere un giudizio ma non \`e esso stesso un oggetto manipolabile --- $\Rightarrow$ vive nel \textbf{metalinguaggio}. \`E una regola, un giudizio, un'affermazione \emph{sul} sistema.

\bigskip

\noindent In altre parole: il linguaggio \`e ci\`o di cui parli; il metalinguaggio \`e ci\`o \emph{con cui} parli. \texttt{zero} \`e una cosa di cui parli (linguaggio); ``\texttt{zero} $:\mathbb{N}$'' \`e una cosa che \emph{dici} --- un'affermazione (metalinguaggio).

\section*{Tre punti dove si inciampa}

\paragraph{(i) I due punti \texttt{:} non sono dentro il linguaggio.}
\texttt{a : A} non \`e un oggetto unico. \texttt{a} e \texttt{A} vivono nel linguaggio; il \texttt{:} che li unisce \`e un'\emph{affermazione sul sistema} (``asserisco che \texttt{a} abita \texttt{A}''). Non si pu\`o passare \texttt{a : A} come argomento a niente: non \`e un oggetto, \`e un giudizio. Il \texttt{:} \`e metalinguaggio puro.

\paragraph{(ii) La barra orizzontale \rule{1cm}{0.4pt} \`e metalinguaggio.}
Separa premesse da conclusione: ``dati questi giudizi sopra, ottieni questo giudizio sotto''. \`E la grammatica con cui si scrivono le regole (Cero~1).

\paragraph{(iii) Il $\vdash$ \`e metalinguaggio, ma fa una cosa diversa dalla barra.}
Sta \emph{dentro una singola premessa}: $x : A \vdash B\,x : \Type$, ``nel contesto in cui assumo $x : A$, vale $B\,x : \Type$''. Costruisce \emph{un} giudizio dicendo sotto quali ipotesi vale (Cero~\ref{cero:contesti}; la grammatica completa di contesti e $\vdash$ \`e in \S\ref{cap:contesti}). La barra combina giudizi interi; il $\vdash$ costruisce un giudizio sotto ipotesi. Non sono lo stesso simbolo.

\bigskip

\noindent Il quarto punto dove si inciampa --- la coppia $\definitorio$ vs.\ $\identita$ --- non lo si pu\`o ancora introdurre per filo perch\'e ci manca il tipo identit\`a; lo chiuderemo in \S\ref{cap:identita}. Per ora basti la promemoria: il $\definitorio$ visto in \S\ref{cap:giudizi} \`e definitorio, vive nel metalinguaggio, non si abita; il $\identita$ che incontreremo in \S\ref{cap:identita} \`e un tipo del linguaggio, si abita con \texttt{refl}.

\section*{Tabella lampo}

\begin{center}
\renewcommand{\arraystretch}{1.15}
\begin{tabular}{lll}
\toprule
\textit{Vedi questo} & \textit{Dov'\`e} & \textit{Perch\'e} \\
\midrule
\texttt{tt}, \texttt{ff}, \texttt{zero}, \texttt{suc n} & linguaggio & termini, manipolabili \\
$\mathbb{B}$, $\mathbb{N}$, \texttt{Vec A n} & linguaggio & tipi, oggetti del sistema \\
\texttt{:} in \texttt{a : A} & meta & giudizio, non oggetto \\
$\definitorio$ definitorio (\S\ref{cap:giudizi}) & meta & asserzione sul sistema \\
\rule{0.9cm}{0.3pt} (la barra) & meta & combina giudizi interi (Cero~1) \\
$\vdash$ & meta & giudizio sotto ipotesi (Cero~\ref{cero:contesti}) \\
coerenza del sistema & meta-meta & non dall'interno --- G\"odel \\
\bottomrule
\end{tabular}
\end{center}

\noindent Le righe sui tipi del linguaggio le completiamo in \S\ref{cap:identita}, quando il tipo identit\`a porter\`a con s\'e altri abitanti del linguaggio.

% ---------------------------------------------------------------------
\chapter{Lo schema F-I-E-C}\label{cap:fiec}

\noindent {\spacedlowsmallcaps{Questa \`e la sezione cardine}}. Ogni tipo in MLTT \`e dato da \textbf{quattro regole}, sempre nello stesso ordine, con la stessa struttura. Non c'\`e altro modo di introdurre un tipo, e non c'\`e altro modo di lavorarci sopra. Tutto quello che si vedr\`a nella Parte~II \`e declinazione di questo schema.

\bigskip

\begin{center}
\renewcommand{\arraystretch}{1.25}
\begin{tabular}{lll}
\toprule
\textit{Regola} & \textit{Cosa dice} & \textit{Dove vive} \\
\midrule
\textbf{Formazione} (F) & quando il tipo esiste & linguaggio \\
\textbf{Introduzione} (I) & come costruisco un abitante (costruttori) & linguaggio \\
\textbf{Eliminazione} (E) & come uso un abitante (eliminatori) & \textbf{metalinguaggio} \\
\textbf{Computazione} (C) & come riduce (regole $\beta$, definitorie) & \textbf{metalinguaggio} \\
\bottomrule
\end{tabular}
\end{center}

\bigskip

Due cose vanno fissate subito, perch\'e ricorrono identiche per ogni tipo.

\section*{Gli eliminatori non sono funzioni}

Non sono funzioni che abitano un tipo: sono \emph{regole} del sistema di deduzione. Quando un testo li presenta come oggetti applicabili ($\indN(\dots)$, $\indB(\dots)$) sta facendo un'\textbf{immersione pedagogica} --- porta qualcosa dal metalinguaggio dentro il linguaggio. L'oggetto cambia natura nel farlo. Da tenere a mente ogni volta che si incontra un eliminatore scritto come se fosse un termine.

\begin{oss}[Test diretto]
\agdainl{suc} si pu\`o passare a una funzione higher-order? S\`i, dunque \`e nel linguaggio. \agdainl{ind-N} in MLTT pura si pu\`o passare? No, \`e una regola del metalinguaggio. Agda --- per ergonomia --- ti far\`a scrivere una funzione di nome \texttt{nat-rec} che si comporta come \agdainl{ind-N}; quella \`e l'immersione pedagogica resa cittadino del linguaggio. Non confondere il gesto pedagogico con la regola.
\end{oss}

\section*{Le regole di computazione sono definitorie}

Non sono teoremi. Non si ha nulla per dimostrarle \emph{dall'interno} del sistema. Si scrivono perch\'e rispecchiano la struttura del tipo. Il sistema funziona perch\'e \`e allineato con quella struttura, non per decreto.

\begin{cero}[Le regole di computazione sono definitorie]\label{cero:comp}
Si assumono, e reggono perch\'e la struttura era gi\`a l\`i.
\end{cero}

\section*{Forma generale di una regola}

Tutte le regole hanno la forma:
\[
\inferrule{\text{premessa}_1 \\ \text{premessa}_2 \\ \dots}{\text{conclusione}}
\]

Premesse e conclusione sono \emph{giudizi} del metalinguaggio (\texttt{a : A}, \texttt{A : Type}, $x : A \vdash B\,x : \Type$, ecc.). La barra dice ``se hai i giudizi sopra, puoi scrivere il giudizio sotto''. Una sequenza di regole applicate fino a un giudizio senza premesse aperte \`e una \emph{derivazione}.

\bigskip

\noindent Da \S\ref{cap:bool} in poi, ogni tipo riceve esattamente queste quattro regole, scritte in questo formato, in quest'ordine.

\section*{Stile di esposizione: alla Little Typer}\label{sec:stile}

\noindent {\spacedlowsmallcaps{Una dichiarazione di metodo}}, prima di entrare in Parte~II.

\bigskip

Esistono due modi storicamente attestati di presentare un tipo induttivo. Il primo --- quello di \emph{The Little Typer} di Friedman e Christiansen, e degli articoli fondativi di Martin-L\"of --- elenca per esteso F-I-E-C e usa \textbf{soltanto} l'eliminatore esplicito ($\indB$, $\indN$, $\Jelim$, $\ldots$) quando arriva il momento di costruire abitanti. Il secondo --- quello di Agda --- scrive una \texttt{data} \texttt{where} per dichiarare il tipo, e poi tutto il resto via \emph{pattern matching}: clausole con costruttori a sinistra del segno uguale, niente eliminatore in vista.

I due stili \emph{non sono equivalenti}: il secondo \`e pi\`u potente, e quanto pi\`u potente lo discuteremo in apertura di Parte~III (\S\ref{cap:prima}). Per ora ci basta la scelta di forma:

\begin{itemize}
\item \textbf{In questo libro le regole le scriviamo alla Little Typer.} Ogni tipo riceve F, I, E, C in forma di regola di inferenza esplicita; l'eliminatore ha un nome ($\indB$, $\indN$, $\ldots$) e si presenta come regola del metalinguaggio.
\item \textbf{Quando arrivano i sorgenti Agda} (Parte~III), useremo il pattern matching come si presenta nei libri di riferimento --- a cominciare da \emph{Verified Functional Programming in Agda} di Stump. Il nostro lavoro di lettura sar\`a sempre: \emph{quale F-I-E-C \`e in atto sotto questa clausola}?
\end{itemize}

\noindent\textbf{Esercizio centrale.} Una volta visti i sei tipi di Parte~II, il lettore dovrebbe essere in grado di guardare una funzione Agda definita per casi e ricostruirne meccanicamente la versione con eliminatore esplicito. La capacit\`a di fare quest'operazione --- almeno nei casi facili --- \`e il segnale che l'F-I-E-C \`e diventata vista, non solo lessico.
% =====================================================================
% PARTE II — I tipi
% =====================================================================
\part{I tipi}

\noindent\itshape Otto capitoli, sempre F-I-E-C.
$\mathbb{B}$ (\S\ref{cap:bool}) per vedere lo schema nel caso pi\`u puro, senza
ricorsione. $\mathbb{N}$ (\S\ref{cap:nat}) per il salto all'induzione vera.
$\Pi$ (\S\ref{cap:pi}) per la decisione fondativa di MLTT, e $\Pi$ in azione
(\S\ref{cap:pi-azione}) per le due passate che la sintassi nasconde. $\forall$
(\S\ref{cap:forall}) come zucchero. $\Sigma$ (\S\ref{cap:sigma}) come duale.
Il tipo identit\`a $\_\identita\_$ (\S\ref{cap:identita}) come sezione cardine
--- l\`i si chiude la distinzione fra le due uguaglianze. $\mathrm{Vec}\,A\,n$
(\S\ref{cap:vec}) come capstone, dove $\mathbb{N}$, $\Pi$ e la dipendenza si
incastrano per la prima volta su un valore numerico.\upshape

\bigskip

% ---------------------------------------------------------------------
\chapter{$\mathbb{B}$ --- il caso pi\`u puro}\label{cap:bool}

Cominciamo da $\mathbb{B}$. \`E il tipo induttivo pi\`u piccolo: due costruttori, zero argomenti, zero ricorsione. L'eliminazione qui \`e \textbf{puro case-split} --- non c'\`e ipotesi induttiva, perch\'e non c'\`e ricorsione su cui ricorrere. Vederlo prima di $\mathbb{N}$ rende esplicito che la ricorsione di $\mathbb{N}$ \`e qualcosa \emph{in pi\`u}, non parte intrinseca di ``eliminare''.

\section*{F --- Formazione}
\[
\inferrule{ }{\mathbb{B} : \Type}
\]
Una regola senza premesse: $\mathbb{B}$ esiste, punto.

\section*{I --- Introduzione}
\[
\inferrule{ }{\tto : \mathbb{B}}
\qquad
\inferrule{ }{\ff : \mathbb{B}}
\]
Due costruttori, entrambi senza premesse.

\section*{E --- Eliminazione (metalinguaggio)}
\[
\inferrule{
b : \mathbb{B} \\
C : \mathbb{B} \to \Type \\
c_t : C\,\tto \\
c_f : C\,\ff
}{
\indB(b, C, c_t, c_f) : C\,b
}
\]

Letto a voce: ``per costruire qualcosa che dipende da un $b : \mathbb{B}$ generico --- cio\`e un abitante di $C\,b$ --- basta darmi il pezzo per $b = \tto$ e il pezzo per $b = \ff$; io metto tutto insieme''. $C$ si chiama \textbf{motivo} dell'eliminazione.

% NUOVO: paragrafo sul motivo, valido per tutti gli eliminatori
\begin{oss}[Il motivo $C$, una volta per tutti gli eliminatori]
Il motivo \`e il concetto pi\`u astratto della Parte~II, e ritorner\`a identico in $\indN$, $\indS$, $\Jelim$, $\indV$. Vale la pena fissarlo subito.

$C$ \`e \textbf{ci\`o che si vuole concludere}, espresso come funzione dal valore eliminato. Scriverlo come $C : \mathbb{B} \to \Type$ significa: ``per ogni $b$ possibile dir\`o cosa intendo costruire''. L'eliminatore promette di consegnare $C\,b$ per il $b$ effettivo; in cambio chiede un pezzo per ogni costruttore.

Due regimi tipici:
\begin{itemize}
\item \textbf{Motivo costante} ($C\,b \definitorio X$ per ogni $b$): si sta solo facendo una funzione $\mathbb{B} \to X$, niente di dipendente. L'eliminatore degenera a un ``if-then-else'' usuale.
\item \textbf{Motivo davvero dipendente} ($C\,b$ varia con $b$): si sta dimostrando una propriet\`a o costruendo un dato il cui \emph{tipo} cambia con il valore. Qui la dipendenza Cero~\ref{cero:dipendenza} viene effettivamente sfruttata.
\end{itemize}
Tutti gli eliminatori che seguono hanno la stessa anatomia: motivo + un caso per costruttore (pi\`u, dove c'\`e ricorsione, un'ipotesi induttiva).
\end{oss}

\section*{C --- Computazione (definitorie, Cero~\ref{cero:comp})}
\begin{align*}
\indB(\tto, C, c_t, c_f) &\definitorio c_t \\
\indB(\ff,  C, c_t, c_f) &\definitorio c_f
\end{align*}

Quando l'argomento \`e un costruttore concreto, l'eliminatore riduce al ramo corrispondente. Nient'altro.

\begin{oss}[Punto da fissare]
Confronta con quanto vedremo per $\mathbb{N}$ in \S\ref{cap:nat}: l\`i $c_t, c_f$ saranno sostituiti da un caso base e un \emph{passo} che riceve un'ipotesi induttiva. Qui no, perch\'e $\tto$ e $\ff$ non contengono altre $\mathbb{B}$. Eliminare = sempre case-split su tutti i costruttori. Indurre = case-split + ricorrere quando il costruttore contiene istanze del tipo stesso. $\mathbb{B}$ ha la prima cosa, non la seconda.
\end{oss}

% ---------------------------------------------------------------------
\chapter{$\mathbb{N}$ --- l'eliminazione diventa induzione}\label{cap:nat}

Stesso schema, ma ora un costruttore (\texttt{suc}) prende un argomento di tipo $\mathbb{N}$. Questo cambia l'eliminazione: nel ramo \texttt{suc k}, oltre a \texttt{k} stesso, riceverai anche $C\,k$ --- il risultato dell'eliminazione \emph{gi\`a calcolato} su \texttt{k}. \`E l'\textbf{ipotesi induttiva}.

\section*{F --- Formazione}
\[
\inferrule{ }{\mathbb{N} : \Type}
\]

\section*{I --- Introduzione}
\[
\inferrule{ }{\zero : \mathbb{N}}
\qquad
\inferrule{n : \mathbb{N}}{\suc\,n : \mathbb{N}}
\]

Due costruttori, ma uno (\texttt{suc}) ha una premessa: prende una $\mathbb{N}$ per produrne un'altra. \`E \textbf{questa} premessa che fa la differenza con $\mathbb{B}$.

\section*{E --- Eliminazione}
\[
\inferrule{
n : \mathbb{N} \\
C : \mathbb{N} \to \Type \\
b : C\,\zero \\
s : (k : \mathbb{N}) \to C\,k \to C\,(\suc\,k)
}{
\indN(n, C, b, s) : C\,n
}
\]

Confronto riga per riga con $\indB$:
\begin{itemize}
% MODIFICATO: uniformato \texttt{tt} -> \tto per coerenza con cap. Bool
\item per $\zero$ serve un valore $b : C\,\zero$ (come $c_t$ per $\tto$);
\item per $\suc\,k$ serve $s : (k : \mathbb{N}) \to C\,k \to C\,(\suc\,k)$ --- non un valore, ma una funzione che prende $k$ \emph{e l'ipotesi induttiva $C\,k$} e produce $C\,(\suc\,k)$.
\end{itemize}

Il $C\,k$ in input a $s$ \`e dove l'eliminazione di $\mathbb{N}$ diventa \textbf{ricorsione primitiva}. \`E letteralmente l'ipotesi induttiva.

\section*{C --- Computazione}
% MODIFICATO: spaziatura uniforme negli argomenti di \indN
\begin{align*}
\indN(\zero,\,C, b, s)    &\definitorio b \\
\indN(\suc\,k,\,C, b, s)  &\definitorio s\,k\,(\indN(k, C, b, s))
\end{align*}

Il caso $\suc\,k$ \`e dove appare la \textbf{ricorsione}: $\indN$ chiama s\'e stesso su $k$, strutturalmente pi\`u piccolo di $\suc\,k$. Questo garantisce terminazione.

% NUOVO: indN in azione --- la somma come esempio concreto di ipotesi induttiva
\section*{$\indN$ in azione --- la somma}

L'addizione \`e l'esempio canonico per vedere l'ipotesi induttiva al lavoro. Definendola tramite $\indN$ con motivo costante $C \definitorio (\lambda\,\_\,.\,\mathbb{N} \to \mathbb{N})$:
\[
n + m \;\definitorio\; \indN\!\bigl(n,\; \lambda\,\_\,.\,\mathbb{N},\; m,\; \lambda\,k\,\mathit{rec}.\,\suc\,\mathit{rec}\bigr)
\]
Letta a voce: ``ricorri su $n$; al caso $\zero$ rispondi $m$; al caso $\suc\,k$ prendi il \emph{risultato gi\`a calcolato} su $k$ --- che chiamo $\mathit{rec}$, cio\`e $k + m$ --- e ne fai $\suc$''. Il parametro $\mathit{rec}$ \`e \textbf{esattamente} l'ipotesi induttiva $C\,k$; il suo essere disponibile \`e ci\`o che distingue $\indN$ dal puro case-split.

In Agda, con pattern matching, lo stesso si scrive nella forma familiare:
\begin{agda}
_+_ : ℕ → ℕ → ℕ
zero  + m = m
suc n + m = suc (n + m)   -- il "(n + m)" qui dentro È l'ipotesi induttiva,
                          --   cioè la chiamata ricorsiva su n
\end{agda}
Le due scritture sono equivalenti: la seconda \`e ci\`o che Agda elabora nella prima.

\section*{Case-split puro \emph{vs} induzione}

Per esplicitare ulteriormente la differenza con $\mathbb{B}$, ecco come sarebbe $\mathbb{N}$ con \emph{solo} case-split (l'analogo di $\indB$):
\[
\inferrule{
n : \mathbb{N} \\
C : \mathbb{N} \to \Type \\
b : C\,\zero \\
s : (k : \mathbb{N}) \to C\,(\suc\,k) \quad \textit{NB: $s$ non riceve $C\,k$}
}{
\mathsf{case\text{-}\mathbb{N}}(n, C, b, s) : C\,n
}
\]

Con questo non si dimostra quasi niente: nel ramo \texttt{suc k} non c'\`e ipotesi su \texttt{k}. La scelta di MLTT \`e di prendere l'\textbf{eliminatore induttivo} invece del solo case-split. Decisione di design, non un'inevitabilit\`a.

% ---------------------------------------------------------------------
\chapter{$\Pi$ --- la dipendenza}\label{cap:pi}

Le funzioni. Generalizzate al punto in cui il \emph{tipo di ritorno} pu\`o dipendere dal \emph{valore} in ingresso.

\section*{F --- Formazione}
\[
\inferrule{A : \Type \\ x : A \vdash B\,x : \Type}{(x : A) \to B\,x : \Type}
\]

$x : A \vdash B\,x : \Type$ significa: \emph{nel contesto in cui ho un $x$ di tipo $A$, $B\,x$ \`e un tipo}. $x$ \`e libera fuori; la regola la chiude dentro il tipo (Cero~\ref{cero:contesti} al lavoro).

Qui c'\`e la \textbf{decisione fondativa} di tutta MLTT:

\begin{cero}[Un termine pu\`o apparire dentro un tipo]\label{cero:dipendenza}
Questa \`e \textsc{la} decisione che definisce MLTT. Prima, hai i tipi semplici, dove $B$ non sa niente di $x$. Dopo, hai i tipi dipendenti.

Non \`e giustificata da niente di precedente. \`E una scelta su \emph{cosa vuoi che il sistema sia}. Potresti scegliere diversamente --- tipi semplici, sistema F, tipi lineari --- e avresti un sistema \textbf{diverso, non sbagliato}. Questo cero distingue MLTT da tutto il resto.
\end{cero}

\section*{I --- Introduzione}
\[
\inferrule{x : A \vdash b\,x : B\,x}{\lambda x.\,b\,x : (x : A) \to B\,x}
\]

Sotto l'ipotesi $x : A$ si costruisce un $b\,x : B\,x$; la regola scarica l'ipotesi producendo la lambda.

\section*{E --- Eliminazione}
\[
\inferrule{f : (x : A) \to B\,x \\ a : A}{f\,a : B\,a}
\]

L'applicazione. \textbf{Banale, nessuna ricorsione} --- $\Pi$ non \`e induttivo, non c'\`e struttura da srotolare in profondit\`a. Contrasto netto con $\indN$.

\section*{C --- Computazione}
\[
(\lambda x.\,b\,x)\,a \definitorio b\,a
\]

La $\beta$-riduzione. Definitoria (Cero~\ref{cero:comp}).

\section*{Nomi alternativi}
\begin{itemize}
\item \textbf{$\Pi$ type} --- nome formale MLTT
\item \textbf{dependent function type} --- descrittivo
\item \textbf{dependent product} --- nome storico, confonde
\end{itemize}

% ---------------------------------------------------------------------
\chapter{$\Pi$ in azione --- i tipi sono termini, ma in due passate}\label{cap:pi-azione}

Sezione di chiarimento sull'uso di $\Pi$. Tre cose vanno fissate, perch\'e Pie (\emph{The Little Typer}) e Agda le fondono sintatticamente.

\section*{I tipi sono termini --- vero, non illusione}

In MLTT esiste un universo $U$ (\agdainl{Set} in Agda) i cui abitanti sono tipi. Quindi una variabile $A : U$ \`e un termine \emph{onesto}: si passa come si passa un numero. \`E il senso vero di ``tipi di prima classe''.

\section*{Ma c'\`e una doppia passata}

Da \emph{The Little Typer}:
\begin{verbatim}
(claim elim-Pair
  (Pi ((A U) (D U) (X U))
      (-> (Pair A D) (-> A D X) X)))

(define elim-Pair
  (lambda (A D X)
    (lambda (p f) (f (car p) (cdr p)))))
\end{verbatim}

Quando il typechecker incontra \verb|(elim-Pair Nat Atom Bool p f)| fa \textbf{due lavori distinti}:

\begin{enumerate}
\item \textbf{Passata 1 --- type checking} (``a compile time''). Contano $A, D, X$. Il checker istanzia $A := \mathrm{Nat}$ ecc.\ nel claim per calcolare di che tipo devono essere \texttt{p} e \texttt{f}.
\item \textbf{Passata 2 --- computazione} (``a runtime/normalizzazione''). Contano \texttt{p} e \texttt{f}. Il corpo \verb|(f (car p) (cdr p))| srotola la coppia. Qui $A, D, X$ sono \textbf{gi\`a stati consumati} nella passata 1 --- non fanno pi\`u nulla.
\end{enumerate}

\textbf{Prima fissi i tipi, poi fai il resto.} L'illusione \emph{non} \`e che i tipi fingano di essere valori --- lo sono. L'illusione \`e che una sola sintassi \verb|(lambda (A D X) (lambda (p f) ...))| nasconde che le due $\lambda$ giocano ruoli temporalmente diversi.

\begin{oss}[Segnale fisico]
$A, D, X$ \emph{non compaiono nel corpo}. Un argomento che il calcolo non usa, ma che il checker pretende, \`e per definizione della passata-tipi.
\end{oss}

\begin{oss}[Conferma teorica --- \emph{type erasure}]
Proprio perch\'e la passata 2 non guarda $A, D, X$, un compilatore pu\`o cancellarli dopo il type checking. Se le passate fossero la stessa cosa, non si potrebbe. Il fatto che si possa \`e la prova che sono separate. (Dove le passate si \emph{toccano} \`e \texttt{Vec A n} (\S\ref{cap:vec}): l\`i un valore \texttt{n} risale dentro un tipo. Si intrecciano, non collassano.)
\end{oss}

\section*{\texttt{elim-Pair} in Agda}

Il claim Pie diventa:

\begin{agda}
-- Pair A D = Sigma A (\ _ -> D)  (prodotto NON dipendente; vedi cap. sigma)

-- versione esplicita: traduzione fedele del claim Pie
elim-Pair : (A : Set) (D : Set) (X : Set) → Pair A D → (A → D → X) → X
elim-Pair A D X p f = f (fst p) (snd p)
--        ^^^^^                          passata 1: tipi, NON usati nel corpo
--              ^^^                       passata 2: valori, USATI nel corpo

-- versione implicita: Agda inferisce A, D, X dal tipo di p
elim-Pair : {A D X : Set} → Pair A D → (A → D → X) → X
elim-Pair p f = f (fst p) (snd p)
\end{agda}

Le graffe \texttt{\{\}} \emph{sono} il riconoscimento sintattico che $A, D, X$ vivono nella passata 1. Chiamando \texttt{elim-Pair miaCoppia miaFunzione}, Agda guarda il tipo di \texttt{miaCoppia} (es. \texttt{Pair Nat Atom}), unifica e riempie $A := \mathrm{Nat}$, $D := \mathrm{Atom}$ da s\'e. La passata 1 avviene, ma invisibile.

\section*{$\Pi$ \`e unario --- gli argomenti multipli sono curry}

Dubbio naturale: \texttt{(A : Set)(D : Set)(X : Set) ->} sembra un $\Pi$ a tre argomenti. \textbf{Non lo \`e.} La regola di formazione di \S\ref{cap:pi} resta \textbf{unaria}. Quella scrittura \`e zucchero per $\Pi$ annidati:

\begin{agda}
(A : Set) → ((D : Set) → ((X : Set) → (Pair A D → ((A → D → X) → X))))
\end{agda}

Un solo $\Pi$ per ogni \texttt{->}, ciascuno con il contesto esteso di un'ipotesi (Cero~\ref{cero:contesti}).

\begin{oss}[Prova che \`e curry e non un $\Pi$ ternario primitivo]
Le applicazioni parziali hanno un tipo.
\begin{agda}
elim-Pair Nat            -- : (D : Set) (X : Set) → Pair Nat D → (Nat → D → X) → X
elim-Pair Nat Atom       -- : (X : Set) → Pair Nat Atom → (Nat → Atom → X) → X
elim-Pair Nat Atom Bool  -- : Pair Nat Atom → (Nat → Atom → Bool) → Bool
\end{agda}

Se fosse un $\Pi$ ternario primitivo, \texttt{elim-Pair Nat} da solo non avrebbe tipo. Il fatto che ce l'ha \`e la dimostrazione che \`e curry --- \textbf{curry dipendente}, perch\'e ogni $\Pi$ vede le variabili legate dai precedenti.
\end{oss}

\bigskip

\noindent\textbf{Da fissare.} \texttt{(A : ..)(B : ..)(C : ..) -> ...} = tre $\Pi$ annidati, non una regola nuova.

% ---------------------------------------------------------------------
\chapter{Il $\forall$ di Agda}\label{cap:forall}

Solo zucchero in pi\`u, ma \`e bene fissarlo perch\'e lo si vede dappertutto.

\begin{agda}
foo : ∀ (x : A) → B x       -- identico a (x : A) → B x
foo : ∀ x → B x             -- identico, ma il tipo di x è inferito da B
\end{agda}

% MODIFICATO: esplicitato che $\forall$ NON è un quantificatore logico primitivo
$\forall$ \textbf{non \`e un quantificatore logico primitivo, n\'e una regola distinta da $\Pi$}: \`e \emph{soltanto zucchero sintattico} per $\Pi$ con il permesso di omettere annotazioni di tipo (quando Agda pu\`o inferirle). Il nome viene dalla lettura logica (``per ogni $x$\ldots''), ma in MLTT non esiste un ``per ogni'' di prima classe separato da $\Pi$: la quantificazione universale \emph{\`e} $\Pi$, via Curry--Howard. Si applica sempre la regola di \S\ref{cap:pi}. Nessuna regola nuova.

% ---------------------------------------------------------------------
\chapter{$\Sigma$ --- il duale di $\Pi$}\label{cap:sigma}

Coppie dipendenti: il secondo componente ha un tipo che dipende dal primo.

\section*{F --- Formazione}
\[
\inferrule{A : \Type \\ x : A \vdash B\,x : \Type}{\Sigma\,x \in A,\,B\,x : \Type}
\]

Stessa premessa di $\Pi$ --- eredita il Cero~\ref{cero:dipendenza}, non ne aggiunge.

\section*{I --- Introduzione}
\[
\inferrule{a : A \\ b : B\,a}{(a, b) : \Sigma\,x \in A,\,B\,x}
\]

\section*{E --- Eliminazione}
\[
\inferrule{
p : \Sigma\,x \in A,\,B\,x \\
C : (\Sigma\,x \in A,\,B\,x) \to \Type \\
f : (a : A) \to (b : B\,a) \to C\,(a, b)
}{
\indS(p, C, f) : C\,p
}
\]

\textbf{Nessuna ricorsione} nel lato destro: una coppia ha esattamente due elementi, non c'\`e struttura da srotolare in profondit\`a. Come $\Pi$, contrasto con $\indN$.

\section*{C --- Computazione}
\[
\indS((a, b), C, f) \definitorio f\,a\,b
\]

\section*{Caso non dipendente: il prodotto}

Quando $B$ non usa $x$, si ottiene il prodotto ordinario:
\[
A \times B \;\;\coloneqq\;\; \Sigma\,A\,(\lambda\,\_\,.\,B)
\]

\texttt{Pair A D} di \S\ref{cap:pi-azione} era questo. Gli accessori \texttt{fst} e \texttt{snd} sono casi particolari di $\indS$ con motivo costante.

\section*{La confusione prodotto/somma}

\begin{center}
\renewcommand{\arraystretch}{1.15}
\begin{tabular}{lll}
\toprule
\textit{Senso} & \textit{Prodotto} & \textit{Somma} \\
\midrule
aritmetico/insiemistico & coppie $A \times B$ & scelta $A + B$ \\
categoriale             & $\Pi$ (funzioni indicizzate) & $\Sigma$ (coppie indicizzate) \\
\bottomrule
\end{tabular}
\end{center}

% MODIFICATO: la confusione "Σ si chiama somma in CT" spostata in nota a pie pagina
\noindent $\Sigma$ generalizza il prodotto \emph{aritmetico} ma si chiama \emph{somma categoriale}\footnote{In teoria delle categorie ``somma'' (o \emph{coprodotto}) \`e il duale di ``prodotto''. Il coprodotto binario \`e $A + B$, e la sua versione indicizzata \`e $\coprod_{x \in A} B\,x$, che \emph{\`e} $\Sigma\,x\!\in\!A,\,B\,x$. Quindi $\Sigma$ \`e ``somma'' nel senso categoriale di coprodotto indicizzato, anche se insiemisticamente i suoi abitanti sono coppie e ``sembra'' un prodotto. Doppia tradizione terminologica, da non confondere.}. Incidente storico, da ignorare. In pratica: $\Pi$ = funzioni dipendenti, $\Sigma$ = coppie dipendenti.

% ---------------------------------------------------------------------
\chapter{$\_\identita\_$ --- il tipo identit\`a}\label{cap:identita}

\epigraph{Una volta capita questa sezione, MLTT smette di sembrare misterioso.}{}

\noindent {\spacedlowsmallcaps{Cardine}}. Qui si chiude il cerchio: anche l'uguaglianza ha le sue quattro regole, \`e un tipo come un altro, e si distingue nettamente dall'uguaglianza definitoria del metalinguaggio.

\bigskip

Fino a \S\ref{cap:giudizi} e \S\ref{cap:test} abbiamo parlato di un'uguaglianza che vive \emph{nel metalinguaggio}: scritta $\equiv$, \`e un giudizio, non un tipo, non si abita, non si dimostra. \`E quella che dice ``queste due cose riducono alla stessa cosa per le regole di computazione''. La chiamiamo \textbf{uguaglianza definitoria}.

Ora introduciamo un \emph{secondo} $\equiv$, completamente diverso, che vive \textbf{nel linguaggio}: \`e un tipo, lo si abita con prove, lo si manipola. Si chiama \textbf{tipo identit\`a} (o \textbf{uguaglianza proposizionale}). Si scrive uguale e questa \`e una sciagura notazionale storica, ma sono due bestie diverse.

\medskip

Per smarcarli scrivo:
\begin{itemize}
\item $a \definitorio b : A$ per l'uguaglianza \textbf{definitoria} (metalinguaggio, \S\ref{cap:giudizi});
\item $a \identita b$ (senza il ``$: \Type$''), o $\mathsf{Id}\,A\,a\,b$ se serve esplicito, per il \textbf{tipo identit\`a} (linguaggio, questa sezione).
\end{itemize}

Quando Agda scrive \agdainl{a ≡ b} intende \textbf{il tipo identit\`a}, sempre. La cosa che gli viene dietro \`e il \emph{secondo} dei due.

\section*{Le quattro regole}

\paragraph{F --- Formazione.}
\[
\inferrule{A : \Type \\ x : A \\ y : A}{x \identita y : \Type}
\]
Dati due elementi di \emph{uno stesso} tipo, si ottiene un tipo: il tipo delle prove che sono uguali.

\paragraph{I --- Introduzione.} (Un solo costruttore --- e questo \`e il punto.)
\[
\inferrule{A : \Type \\ x : A}{\refl : x \identita x}
\]

C'\`e \textbf{un solo} costruttore, $\refl$, e abita \textbf{soltanto} $x \identita x$. Cio\`e: l'unico modo canonico di costruire una prova di uguaglianza \`e quando i due lati sono \textbf{gi\`a la stessa cosa}.

\bigskip

\noindent\textbf{Qui salta fuori la cosa cruciale}: ``gi\`a la stessa cosa'' \emph{in che senso}? Nel senso dell'\textbf{uguaglianza definitoria del metalinguaggio}, cio\`e dopo aver applicato tutte le regole di computazione viste finora ($\indN$ su \texttt{zero}, $\beta$-riduzione di $\Pi$, ecc.).

\begin{regola}[Il ponte]\label{reg:ponte}
$\refl$ ha tipo $x \identita y$ (linguaggio) \textbf{se e solo se} $x \definitorio y : A$ (definitorio, metalinguaggio).
\end{regola}

$\refl$ \emph{colma la trincea} tra meta e linguaggio. \`E il termine che dice ``i due lati riducono allo stesso normale, e io ti porto questa informazione dentro il sistema sotto forma di prova''.

Quindi le due uguaglianze non sono indipendenti: il tipo identit\`a \emph{si appoggia} alla definitoria. La definitoria sta ``sotto'' e fa il lavoro silenzioso di riduzione; il tipo identit\`a sta ``sopra'' e ne raccoglie il risultato come prova manipolabile.

% NUOVO: esempi Agda di refl che funziona e refl che fallisce
\bigskip

\noindent\textbf{In pratica, con $+$ definita come in \S\ref{cap:nat}}:
\begin{agda}
-- (1) refl funziona: entrambi i lati riducono a 4
_ : 2 + 2 ≡ 4
_ = refl

-- (2) refl funziona ANCHE qui: 2+2 e 3+1 hanno lo stesso normale (= 4)
_ : 2 + 2 ≡ 3 + 1
_ = refl

-- (3) refl funziona: zero + n riduce a n per la prima equazione di +
_ : ∀ {n} → zero + n ≡ n
_ = refl

-- (4) refl FALLISCE: n + zero NON riduce, perché + ricorre sul PRIMO
--     argomento e n è una variabile (nessuna delle equazioni si attiva).
--     Definitoriamente non è uguale a n; serve induzione su n per dimostrarlo.
-- _ : ∀ {n} → n + zero ≡ n
-- _ = refl   -- errore di typecheck: n + zero != n
\end{agda}

\noindent (1)--(3) sono ``stessa normale dopo riduzione'', quindi la Regola~\ref{reg:ponte} li accetta. (4) non lo \`e: \emph{semanticamente} $n + \zero$ vale $n$, ma \emph{definitoriamente} non riduce, perch\'e la regola pertinente di $+$ \`e bloccata sulla variabile $n$. Serve $\indN$ per dimostrarlo --- e il risultato non sar\`a $\refl$, ma un termine costruito induttivamente.

\medskip

\noindent\textbf{Da fissare.} ``$\refl$ funziona'' $\iff$ ``i due lati riducono alla stessa forma normale per le regole di computazione''. Tutto il resto richiede prove genuine.

\paragraph{E --- Eliminazione.}
\[
\inferrule{
p : x \identita y \\
C : (y : A) \to (x \identita y) \to \Type \\
c : C\,x\,\refl
}{
\Jelim(p, C, c) : C\,y\,p
}
\]

L'eliminatore $\Jelim$ (\emph{path induction}). Per costruire qualcosa che dipende da una prova generica $p : x \identita y$, basta saperlo costruire nel caso canonico $p = \refl$ (dove necessariamente $y = x$).

\paragraph{C --- Computazione.}
\[
\Jelim(\refl, C, c) \definitorio c
\]
Quando la prova \`e canonica, $\Jelim$ restituisce il caso base.

% NUOVO: subst come corollario immediato di J --- J in azione
\subsection*{$\Jelim$ in azione --- $\mathrm{subst}$ come corollario immediato}

Il principio della \emph{path induction} \`e formidabile ma astratto. Il suo primo corollario \`e una funzione che usiamo continuamente: \textbf{trasporto} di una propriet\`a lungo un'uguaglianza.
\[
\mathrm{subst} : \{A : \Type\}\,(P : A \to \Type)\,\{x\,y : A\} \to x \identita y \to P\,x \to P\,y
\]
Letta a voce: ``se $x \identita y$ e ho un abitante di $P\,x$, ne ottengo uno di $P\,y$''. Si ottiene da $\Jelim$ scegliendo come motivo $C\,y\,p \definitorio P\,y$ (cio\`e: il motivo dipende solo dall'estremo, ignora la prova):
\begin{agda}
subst : {A : Set} (P : A → Set) {x y : A} → x ≡ y → P x → P y
subst P p px = J p (λ y _ → P y) px
--                  ^^^^^^^^^^^
--                  motivo: "non mi importa della prova,
--                   voglio dire qualcosa su y"
\end{agda}
\noindent Il caso base $C\,x\,\refl \definitorio P\,x$ \`e gi\`a dato (lo \`e \texttt{px}); la computazione di $\Jelim$ su $\refl$ restituisce esattamente quel \texttt{px}. Quando invece la prova non \`e $\refl$ ma una qualunque, $\Jelim$ ``trasporta'' il dato lungo l'identit\`a.

\begin{oss}[Perch\'e questo importa]
$\mathrm{subst}$ \`e ci\`o che permette il \emph{rewriting} nelle dimostrazioni in Agda (la parola chiave \texttt{rewrite} si elabora in chiamate a $\mathrm{subst}$). Senza $\Jelim$, non ci sarebbe nessun modo di ``usare'' una prova di uguaglianza per cambiare il tipo di un termine. $\Jelim$ \`e l'astratto; $\mathrm{subst}$ \`e il primo concreto.
\end{oss}

\section*{Ricapitolazione della distinzione}

\begin{center}
\renewcommand{\arraystretch}{1.2}
\begin{tabular}{lll}
\toprule
 & \textit{Uguaglianza definitoria} & \textit{Tipo identit\`a} \\
\midrule
Dove vive & \textbf{metalinguaggio} & \textbf{linguaggio} \\
Forma     & giudizio $a \definitorio b : A$ & tipo $a \identita b$ \\
Si abita? & no, \`e un'affermazione & s\`i, con $\refl$ (e via $\Jelim$) \\
Si scrive? & nelle regole, non in Agda & in Agda: \agdainl{a ≡ b} \\
Viene da  & regole di computazione (Cero~\ref{cero:comp}) & regole F-I-E-C di \S\ref{cap:identita} \\
Ponte     & \multicolumn{2}{c}{abitato da $\refl$ $\iff$ vale la definitoria (Regola~\ref{reg:ponte})} \\
\bottomrule
\end{tabular}
\end{center}

\bigskip

\noindent Quindi, quando in Agda si scrive:
\begin{agda}
qualcosa : x ≡ y
qualcosa = refl
\end{agda}
si \emph{abita} un tipo (linguaggio); Agda accetta solo se \texttt{x} e \texttt{y} sono definitoriamente uguali (metalinguaggio). I due livelli parlano tramite $\refl$. \textbf{Questa \`e la cosa pi\`u importante da tenere chiara di tutto MLTT.}

\section*{Quarto punto dove si inciampa}

In \S\ref{cap:test} avevamo promesso un quarto punto, da chiudere appena introdotto il tipo identit\`a. Eccolo: i due $\equiv$ si scrivono uguale ma sono due cose diverse. Sopra c'\`e la regola per distinguerli --- guarda \emph{cosa} stai dicendo: un'affermazione sul sistema, o un tipo che vuoi abitare?

\medskip

Completiamo la tabella lampo di \S\ref{cap:test}:

\begin{center}
\renewcommand{\arraystretch}{1.15}
\begin{tabular}{lll}
\toprule
\textit{Vedi questo} & \textit{Dov'\`e} & \textit{Perch\'e} \\
\midrule
$\equiv$ in $(\lambda x.\,b)\,a \equiv b\,a$ (regola di computazione) & meta & uguaglianza definitoria, asserzione \\
$\equiv$ in $x \equiv y : \Type$ (Agda: \texttt{\_$\equiv$\_}) & linguaggio & tipo identit\`a, si abita con $\refl$ \\
$\refl$ & linguaggio & termine, abita il tipo identit\`a \\
$\indB$, $\indN$, $\indS$, $\Jelim$ & meta & eliminatori = regole di inferenza \\
\bottomrule
\end{tabular}
\end{center}

% ---------------------------------------------------------------------
\chapter{$\mathrm{Vec}\,A\,n$ --- tutto messo insieme}\label{cap:vec}

Capstone della Parte~II: primo tipo dove $\mathbb{N}$, $\Pi$ e la dipendenza Cero~\ref{cero:dipendenza} si incastrano davvero. \texttt{Vec A n} \`e il tipo delle liste di elementi di $A$ con \emph{esattamente} $n$ elementi. \textbf{$n$ vive nel tipo} --- il typechecker conosce la lunghezza a compile time. Questo \`e il Cero~\ref{cero:dipendenza} al lavoro su un valore numerico (\`e anche il punto dove le due passate di \S\ref{cap:pi-azione} si \emph{toccano}: $n$ \`e un valore che entra in un tipo).

\section*{F --- Formazione}
\[
\inferrule{A : \Type \\ n : \mathbb{N}}{\mathrm{Vec}\,A\,n : \Type}
\]

Il tipo dipende da un \emph{valore} $n : \mathbb{N}$. Non \`e \texttt{Vec A} genericamente: \`e un tipo diverso per ogni lunghezza.

\section*{I --- Introduzione}
\[
\inferrule{ }{\nil : \mathrm{Vec}\,A\,\zero}
\qquad
\inferrule{a : A \\ as : \mathrm{Vec}\,A\,n}{\cons\,a\,as : \mathrm{Vec}\,A\,(\suc\,n)}
\]

\texttt{cons} incrementa $n$ nel tipo risultante. Non per calcolo a runtime, \textbf{per costruzione}.

\section*{E --- Eliminazione (metalinguaggio)}
\[
\inferrule{
vs : \mathrm{Vec}\,A\,n \\
C : (n : \mathbb{N}) \to \mathrm{Vec}\,A\,n \to \Type \\
b : C\,\zero\,\nil \\
s : (k : \mathbb{N}) \to (a : A) \to (as : \mathrm{Vec}\,A\,k) \to C\,k\,as \to C\,(\suc\,k)\,(\cons\,a\,as)
}{
\indV(vs, C, b, s) : C\,n\,vs
}
\]

Stessa struttura di $\indN$: caso base, e passo che riceve l'ipotesi induttiva $C\,k\,as$. Il motivo $C$ dipende sia dalla lunghezza che dal vettore.

\section*{C --- Computazione}
\begin{align*}
\indV(\nil,         C, b, s) &\definitorio b \\
\indV(\cons\,a\,as, C, b, s) &\definitorio s\,k\,a\,as\,(\indV(as, C, b, s))
\end{align*}

Ricorsione su \texttt{as}, strutturalmente pi\`u corto di \texttt{cons a as}. Stesso schema di $\indN$.

\section*{Nota sintattica: parametri \emph{vs} indici}

\begin{agda}
data Vec (A : Set) : ℕ → Set where
--       ^^^^^^^^   ^^^^^^^^^
--       parametro  indice
\end{agda}

\begin{itemize}
\item \textbf{Parametro} \texttt{(A : Set)}: prima dei \texttt{:}, fisso per tutti i costruttori.
\item \textbf{Indice} \texttt{: $\mathbb{N}$ -> Set}: dopo i \texttt{:}, pu\`o cambiare costruttore per costruttore (\texttt{nil} lo fissa a \texttt{zero}, \texttt{cons} lo porta da $n$ a $\suc\,n$).
\end{itemize}

Se $n$ fosse un parametro, \texttt{cons} non potrebbe cambiarlo. Questa \`e una decisione di design (\S\ref{cap:ceri}).

% NUOVO: avvertimento sull'inferenza fragile per i tipi con indici
\section*{Avvertimento --- inferenza fragile sugli indici}

L'inferenza di tipo di Agda \`e completa per i \emph{parametri}, ma \textbf{non lo \`e per gli indici}. Quando un'operazione richiede che due indici risultino definitoriamente uguali e non lo sono per pura riduzione, l'unificatore si arresta e devi fornire una prova di uguaglianza (tipicamente con \texttt{rewrite} o $\mathrm{subst}$, \S\ref{cap:identita}). Tre situazioni tipiche:

\begin{enumerate}
\item \textbf{Concatenazione e lemmi di associativit\`a/commutativit\`a sugli indici}:
\begin{agda}
_++_ : ∀ {A m n} → Vec A m → Vec A n → Vec A (m + n)

-- e ora, per dimostrare l'associatività di ++ a livello di vettori,
-- serve PRIMA un lemma su ℕ (associatività di +) e POI un rewrite,
-- perché (m + n) + p non è DEFINITORIAMENTE uguale a m + (n + p):
-- lo è solo proposizionalmente, e Agda non può forzarlo da solo.
\end{agda}

\item \textbf{Funzioni con indice non-canonico in posizione di pattern}:
\begin{agda}
head : ∀ {A n} → Vec A (suc n) → A
head (cons a _) = a
-- ok: l'indice è "suc n", l'unificazione vede subito che il vettore
-- non può essere nil (che vorrebbe l'indice zero). Bene.

-- Ma se chiami head su un Vec A m con m generico, prima devi
-- convincere Agda che m ≡ suc k per qualche k.
\end{agda}

\item \textbf{Funzioni la cui forma del risultato non si decide per riduzione}: capita spesso quando l'indice contiene una funzione bloccata su una variabile (come $n + \zero$ rispetto a $n$).
\end{enumerate}

\noindent\textbf{Da fissare.} Tipi dipendenti su indici sono potenti, ma il prezzo \`e che l'inferenza non \`e ``magica'': dove la riduzione si ferma, devi fornire tu il ponte con una prova proposizionale. Questo non \`e un difetto di Agda; \`e la conseguenza diretta della Regola~\ref{reg:ponte} --- definitorio sta sotto, proposizionale sopra.

% NUOVO: tabella riassuntiva finale di tutta la Parte II
\section*{Quadro d'insieme --- la Parte~II in una tabella}

I sei tipi della Parte~II, allineati sulle stesse colonne. Stessa anatomia F-I-E-C; ci\`o che cambia \`e la presenza (o l'assenza) di ricorsione, e da cosa dipende il motivo $C$.

\begin{center}
\renewcommand{\arraystretch}{1.25}
\small
\begin{tabular}{llllp{4.6cm}}
\toprule
\textit{Tipo} & \textit{F-I-E-C} & \textit{Ricorsivo} & \textit{Motivo dipende da} & \textit{Uso principale} \\
\midrule
$\mathbb{B}$ & s\`i & no & $b : \mathbb{B}$ & case-split puro; pi\`u piccolo induttivo possibile \\
$\mathbb{N}$ & s\`i & s\`i & $n : \mathbb{N}$ & induzione e ricorsione primitiva \\
$\Pi$        & s\`i & no & --- (non induttivo) & funzioni dipendenti; decisione fondativa MLTT \\
$\Sigma$     & s\`i & no & $p : \Sigma\,x\!\in\!A,\,B\,x$ & coppie dipendenti; duale di $\Pi$ \\
$\_\identita\_$ & s\`i & no & $y : A$ \emph{e} $p : x \identita y$ & uguaglianza interna; ponte meta--linguaggio \\
$\mathrm{Vec}\,A\,n$ & s\`i & s\`i & $n : \mathbb{N}$ \emph{e} $vs : \mathrm{Vec}\,A\,n$ & dipendenza su un valore numerico; capstone \\
\bottomrule
\end{tabular}
\end{center}

\medskip

\noindent\textbf{Letture trasversali.}
\begin{itemize}
\item \emph{Ricorsione = costruttore che cita il tipo stesso.} $\mathbb{B}$ e $\Sigma$ non hanno tale costruttore; $\mathbb{N}$ ($\suc$) e $\mathrm{Vec}$ ($\cons$) s\`i. $\Pi$ e $\_\identita\_$ non sono induttivi in senso stretto.
\item \emph{Motivo a una variabile vs.\ a due.} $\indB$, $\indN$, $\indS$ hanno motivi su una variabile (il valore eliminato). $\Jelim$ e $\indV$ ne hanno due, perch\'e il tipo dipende da pi\`u di un dato (in $\Jelim$: l'altro estremo e la prova; in $\mathrm{Vec}$: la lunghezza e il vettore).
\item \emph{Dipendenza Cero~\ref{cero:dipendenza}.} $\Pi$ la introduce; $\Sigma$ la eredita; $\_\identita\_$ e $\mathrm{Vec}$ la sfruttano in modo essenziale (un \emph{valore} entra in un \emph{tipo}). $\mathbb{B}$ e $\mathbb{N}$ no, di per s\'e: la usano solo quando si scelgono motivi non costanti.
\end{itemize}

\noindent Da qui in poi (Parte~III) tutti gli ingredienti sono pronti: le dimostrazioni saranno termini in tipi costruiti con queste sei regole F-I-E-C, niente di pi\`u.
% =====================================================================
% PARTE III — Lavorare nel sistema
% =====================================================================
\part{Lavorare nel sistema}

\noindent\itshape Quattro capitoli che chiudono la traversata. Prima un
preambolo (\S\ref{cap:prima}) per dichiarare cosa cambia entrando in
Parte~III: i sorgenti Agda useranno pattern matching, e questo ha un
costo che va esplicitato (quinto cero). Poi la logica non importata ma
riconosciuta dentro F-I-E-C, fino a osservare Curry-Howard come gesto
finale (\S\ref{cap:logica}). Un kit di costruzioni canoniche --- i
\emph{tipi notevoli} (\S\ref{cap:notevoli}) --- che fa da palestra
prima delle prove vere. E infine l'esercizio di Stump sulla
commutativit\`a di \texttt{\&\&} letto come tre F-I-E-C intrecciati
(\S\ref{cap:and-comm}), apertura del compendio al capitolo~1 del libro
in oggetto di studio.\upshape

\bigskip

% ---------------------------------------------------------------------
\chapter{Prima del compendio}\label{cap:prima}

\noindent {\spacedlowsmallcaps{Cosa cambia, entrando in Parte~III}}. In Parte~II si scrivevano regole. Da qui in poi si \emph{leggono} sorgenti, ed \`e onesto dichiarare in apertura che cosa si sta facendo quando una clausola Agda viene tradotta in F-I-E-C.

\bigskip

\noindent\textbf{La situazione.} \emph{Verified Functional Programming in Agda} (e qualsiasi sorgente Agda reale) usa il \emph{pattern matching} come modo standard di definire funzioni e dimostrazioni. Una clausola del tipo

\begin{agda}
f : (b : 𝔹) → C b
f tt = ...
f ff = ...
\end{agda}

non scrive $\indB$ da nessuna parte. Il nostro compito di lettura sar\`a riconoscere, sotto ogni clausola di questo genere, il F-I-E-C in atto: che cosa \`e l'eliminazione, che cosa la computazione, che cosa l'introduzione del risultato.

Nel caso semplice di $\mathbb{B}$ la traduzione \`e meccanica. Su $\mathbb{N}$ e $\mathrm{Vec}$ comincer\`a a non esserlo: il pattern matching dipendente \emph{eccede} ci\`o che $\indN$, $\indV$ ecc.\ riescono a fare da soli. Va detto chiaramente, ed \`e il contenuto del quinto cero.

\section*{Il quinto cero}

\begin{cero}[Il pattern matching dipendente eccede MLTT pura]\label{cero:pm}
Il pattern matching che vediamo in Agda non \`e zucchero per gli eliminatori di Parte~II. Nella sua forma generale o \textbf{assume l'assioma K} (\emph{uniqueness of identity proofs}, ogni prova di $x \identita x$ \`e definitoriamente $\refl$), oppure --- in modalit\`a \texttt{--without-K} --- si appoggia a un'\textbf{elaborazione non banale} (Goguen-McBride-McKinna, Cockx-Abel) basata su unificazione di indici, \emph{dot patterns}, \emph{absurd patterns}, \emph{with-abstraction}.

In entrambi i casi: ci\`o che scriviamo come clausole con costruttori a sinistra \emph{non \`e} immediatamente un'applicazione di $\indB$, $\indN$, $\Jelim$ ecc. C'\`e un compilatore in mezzo. Quando in questo libro leggiamo una clausola come ``tre F-I-E-C intrecciati'' (\S\ref{cap:and-comm}) stiamo facendo un'operazione di traduzione che, fuori dai casi pi\`u semplici, non \`e meccanica.
\end{cero}

\section*{Dove morde la profondit\`a: un esempio breve}

L'esempio canonico \`e il matching di $\refl$ contro una prova generica $p : x \identita y$. Vedremo la situazione completa in \S\ref{cap:notevoli} con $\mathsf{sym}$ e $\mathsf{trans}$, ma il punto \`e fissabile in poche righe.

Scrivere

\begin{agda}
sym : {A : Set} {x y : A} → x ≡ y → y ≡ x
sym refl = refl
\end{agda}

\noindent \emph{sembra} dire: ``c'\`e un solo caso, quello in cui la prova \`e $\refl$, e nel ramo $\refl$ il sistema sa che $y = x$''. Ma da dove viene quel ``il sistema sa che $y = x$''? In MLTT pura, l'unico modo di estrarre informazione da $p : x \identita y$ \`e $\Jelim$ (\S\ref{cap:identita}): si specifica il motivo $C : (y : A) \to (x \identita y) \to \Type$ e si abita il caso $C\,x\,\refl$.

\medskip

\textbf{Il matching di Agda fa di pi\`u}: nel ramo $\refl$ \emph{istanzia} $y := x$ nel resto del contesto. Quell'istanziazione, in piena generalit\`a, \`e l'assioma K. In \texttt{--without-K} si controlla che l'istanziazione sia legittima sul piano dell'unificazione --- e questa verifica \`e un algoritmo, non una regola di MLTT.

\begin{oss}[Conseguenza pratica]
Su tipi senza uguaglianze proposizionali non banali fra costruttori --- $\mathbb{B}$, $\mathbb{N}$ per gli esempi facili --- la traduzione clausola$\to$eliminatore \`e meccanica e non chiama in causa K. \texttt{\&\&-comm} (\S\ref{cap:and-comm}) \`e proprio uno di questi casi. Quando lavoreremo su prove che \emph{matchano} su $\refl$ in modo essenziale (\S\ref{cap:notevoli} per $\mathsf{sym}$, $\mathsf{trans}$, $\mathsf{cong}$, $\mathsf{subst}$) il cero~\ref{cero:pm} sar\`a in scena, e va saputo.
\end{oss}

\section*{Patto di lettura per Parte~III}

\noindent\textbf{(i)} Ogni clausola Agda di Parte~III verr\`a accompagnata, almeno la prima volta che il pattern compare, dalla sua traduzione esplicita in F-I-E-C (eliminatore + motivo + casi base/passo).

\noindent\textbf{(ii)} Quando la traduzione \`e \emph{meccanica} (nessun matching essenziale su $\refl$, nessun uso di dot pattern o absurd pattern non banale) ce lo segneremo come ``elaborazione triviale''.

\noindent\textbf{(iii)} Quando \`e \emph{non meccanica} il cero~\ref{cero:pm} sar\`a in scena, lo dichiareremo, e indicheremo a quale lato (K o \texttt{--without-K}) ci stiamo appoggiando.

\bigskip

\noindent Con questo patto in tasca, possiamo entrare in Parte~III sapendo cosa stiamo importando e cosa no.

% ---------------------------------------------------------------------
\chapter{La logica che affiora}\label{cap:logica}

\noindent {\spacedlowsmallcaps{Nessun cero nuovo qui}}. Sono \textbf{osservazioni}, non assunzioni. La logica non \`e importata: \`e gi\`a l\`i nella struttura F-I-E-C.

\bigskip

\noindent\textbf{Gesto del capitolo.} Non diremo ``ora decidiamo che i tipi sono proposizioni e i termini sono prove''. Quella sarebbe un'assunzione --- un cero. Diremo invece: \emph{guardiamo le regole F-I-E-C dei tipi che abbiamo costruito in Parte~II, e notiamo che hanno la forma delle regole della logica intuizionista del primo ordine}. La differenza fra le due mosse \`e tutto il capitolo. La prima impone una lettura; la seconda riconosce una struttura. Il nome di questa coincidenza --- \emph{corrispondenza di Curry-Howard} --- arriva alla fine come osservazione su quanto avr\`a gi\`a parlato da s\'e.

\section*{Implicazione: $\Pi$ non dipendente}

Quando $B$ non dipende da $x$, la regola di formazione di $\Pi$ (\S\ref{cap:pi}) si specializza:
\[
\inferrule{A : \Type \\ B : \Type}{A \to B : \Type}
\]

\noindent\textbf{Lettura logica.} Premesse: $A$ \`e una proposizione, $B$ \`e una proposizione. Conclusione: $A \to B$ \`e una proposizione --- l'implicazione $A \Rightarrow B$.

Le altre due regole di $\Pi$ si leggono altrettanto direttamente:
\begin{itemize}
\item Introduzione $\lambda x.\,b : A \to B$ sotto l'ipotesi $x : A$: \textbf{dimostrazione per ipotesi}. Si assume $A$, si deriva $B$, si chiude scaricando l'ipotesi. \`E la regola $\Rightarrow$-introduzione della deduzione naturale.
\item Eliminazione $f\,a : B$ dati $f : A \to B$ e $a : A$: \textbf{modus ponens}. La regola $\Rightarrow$-eliminazione.
\end{itemize}

\noindent Niente \`e stato deciso. Lo scaricamento delle ipotesi --- gesto centrale della deduzione naturale --- \emph{coincide con} l'introduzione di $\Pi$, che era stata scritta per ragioni di calcolo (descrivere come si costruisce una funzione).

\section*{Congiunzione: $\Sigma$ non dipendente}

Quando $B$ non dipende da $x$, la formazione di $\Sigma$ (\S\ref{cap:sigma}) d\`a:
\[
\inferrule{A : \Type \\ B : \Type}{A \times B : \Type}
\]

\noindent\textbf{Lettura logica.} $A \times B$ \`e la proposizione $A \wedge B$.

\begin{itemize}
\item Introduzione $(a, b) : A \times B$ da $a : A$ e $b : B$: \textbf{regola $\wedge$-introduzione}.
\item Eliminazioni $\mathsf{fst}$, $\mathsf{snd}$ (caso particolare di $\indS$ con motivo costante): \textbf{regole $\wedge$-eliminazione 1 e 2}.
\end{itemize}

\section*{Falsit\`a: $\bot$}

Il tipo vuoto in Agda:

\begin{agda}
data ⊥ : Set where
  -- nessun costruttore
\end{agda}

Le sue quattro regole F-I-E-C esplicitate:

\[
\inferrule{ }{\bot : \Type}
\qquad
\textit{(nessuna I)}
\qquad
\inferrule{e : \bot \\ C : \bot \to \Type}{\mathsf{ind}\text{-}\bot(e, C) : C\,e}
\qquad
\textit{(nessuna C)}
\]

Una Formazione, \emph{zero} Introduzioni, un'Eliminazione, \emph{zero} Computazioni. Eliminatore di $\bot$ in Agda:

\begin{agda}
⊥-elim : {A : Set} → ⊥ → A
⊥-elim ()
\end{agda}

L'\emph{absurd pattern} \texttt{()} \`e la realizzazione sintattica di ``zero costruttori, niente da matchare, l'esaustivit\`a si chiude da sola''.

\noindent\textbf{Lettura logica.}
\begin{itemize}
\item $\bot$ non \`e abitato (nessuna introduzione): \textbf{$\bot$ \`e la falsit\`a}.
\item L'eliminatore d\`a $\bot \to A$ per ogni $A$: \textbf{ex falso quodlibet}.
\end{itemize}

\section*{Negazione: una definizione, non una primitiva}

Non c'\`e una regola F-I-E-C per la negazione. \`E una \emph{definizione}:
\[
\neg A \;\;\coloneqq\;\; A \to \bot
\]

\noindent ``Negare $A$'' significa ``mostrare che da $A$ si deriverebbe l'assurdo''. Lettura intuizionista canonica della negazione, ottenuta gratis dall'avere $\to$ e $\bot$. Una prova di $\neg A$ \`e una funzione che riceve un ipotetico abitante di $A$ e produce un abitante di $\bot$ (cosa impossibile in concreto: \`e la dimostrazione di impossibilit\`a).

\section*{Quantificatore universale: $\Pi$ dipendente}

Senza specializzare, la regola di formazione di $\Pi$ (\S\ref{cap:pi}) \`e gi\`a:
\[
\inferrule{A : \Type \\ x : A \vdash B\,x : \Type}{(x : A) \to B\,x : \Type}
\]

\noindent\textbf{Lettura logica.} ``Per ogni $x$ di tipo $A$, vale $B\,x$''. \`E $\forall x : A.\,B\,x$.

L'introduzione $\lambda x.\,b\,x$ \`e: \emph{costruisci la prova di $B\,x$ uniformemente in $x$, e ottieni la prova del $\forall$}. L'eliminazione $f\,a : B\,a$ \`e: \emph{specializza il $\forall$ all'elemento $a$}.

\noindent Stessa regola di $\Pi$, due letture. Non ne abbiamo aggiunta nessuna nuova.

\section*{Quantificatore esistenziale: $\Sigma$ dipendente}

Stesso gesto sull'altra regola (\S\ref{cap:sigma}):
\[
\inferrule{A : \Type \\ x : A \vdash B\,x : \Type}{\Sigma\,x \in A,\,B\,x : \Type}
\]

\noindent\textbf{Lettura logica.} ``Esiste $x$ di tipo $A$ tale che $B\,x$''. \`E $\exists x : A.\,B\,x$.

L'introduzione $(a, b) : \Sigma\,x \in A,\,B\,x$ richiede un \emph{testimone} $a$ e una \emph{prova} $b : B\,a$. Questa \`e una propriet\`a forte della logica intuizionista: l'esistenza \`e \textbf{costruttiva}, una prova di $\exists x.\,P\,x$ \emph{esibisce} l'$x$ in questione. Non si pu\`o dimostrare ``esiste qualcosa con propriet\`a $P$'' senza dire \emph{cosa} ce l'ha.

\begin{oss}[Confronto con la matematica classica]
In matematica classica si dimostra spesso ``esiste $x$ tale che $\ldots$'' assumendo per assurdo la negazione e derivando una contraddizione. La prova esiste, il testimone no. In MLTT questa strategia \emph{non funziona}: per abitare $\Sigma\,x \in A,\,B\,x$ devi produrre $(a, b)$. \`E lo stesso fenomeno che vedremo a fine capitolo per la doppia negazione e LEM.
\end{oss}

\section*{Disgiunzione: il coproduct $A \uplus B$}\label{sec:coproduct}

Qui c'\`e una connettiva che ci manca. Non l'abbiamo introdotta in Parte~II perch\'e non serviva per la traversata dei sei tipi fondativi. La introduciamo adesso, perch\'e la lettura logica la chiama: \emph{disgiunzione costruttiva}.

\medskip

\noindent F-I-E-C completa del coproduct (o \emph{somma disgiunta}):

\paragraph{F --- Formazione.}
\[
\inferrule{A : \Type \\ B : \Type}{A \uplus B : \Type}
\]

\paragraph{I --- Introduzione.} (Due costruttori.)
\[
\inferrule{a : A}{\mathsf{inj}_1\,a : A \uplus B}
\qquad
\inferrule{b : B}{\mathsf{inj}_2\,b : A \uplus B}
\]

\paragraph{E --- Eliminazione.}
\[
\inferrule{
e : A \uplus B \\
C : (A \uplus B) \to \Type \\
f : (a : A) \to C\,(\mathsf{inj}_1\,a) \\
g : (b : B) \to C\,(\mathsf{inj}_2\,b)
}{
\mathsf{ind}\text{-}\uplus(e, C, f, g) : C\,e
}
\]

\paragraph{C --- Computazione.}
\begin{align*}
\mathsf{ind}\text{-}\uplus(\mathsf{inj}_1\,a, C, f, g) &\definitorio f\,a \\
\mathsf{ind}\text{-}\uplus(\mathsf{inj}_2\,b, C, f, g) &\definitorio g\,b
\end{align*}

Stesso schema visto in Parte~II: due costruttori, nessuna ricorsione (perch\'e n\'e $\mathsf{inj}_1$ n\'e $\mathsf{inj}_2$ contengono istanze di $A \uplus B$), eliminazione per casi.

In Agda:

\begin{agda}
data _⊎_ (A B : Set) : Set where
  inj₁ : A → A ⊎ B
  inj₂ : B → A ⊎ B
\end{agda}

\noindent\textbf{Lettura logica.} $A \uplus B$ \`e la proposizione $A \vee B$.

\begin{itemize}
\item Introduzioni: regole $\vee$-introduzione 1 e 2.
\item Eliminazione: regola di \emph{ragionamento per casi} ($\vee$-eliminazione). Per concludere $C$ da $A \vee B$, mostra come concludere $C$ assumendo $A$, e come concludere $C$ assumendo $B$.
\end{itemize}

\textbf{Costruttivit\`a della disgiunzione.} Una prova di $A \uplus B$ ti dice \emph{quale dei due} vale. Non c'\`e modo di abitarlo senza decidere fra $\mathsf{inj}_1$ e $\mathsf{inj}_2$. \`E lo stesso fenomeno dell'esistenziale: la disgiunzione \`e \emph{informativa}.

\section*{Uguaglianza}

Il tipo identit\`a $\_\identita\_$ (\S\ref{cap:identita}) \`e gi\`a la lettura logica dell'uguaglianza fra termini di uno stesso tipo. $\refl$ \`e l'unica introduzione, e via $\Jelim$ si dimostrano le propriet\`a usuali (riflessivit\`a per costruzione, simmetria, transitivit\`a, sostituzione --- le faremo in \S\ref{cap:notevoli}). Niente da aggiungere qui, salvo notare che \textbf{non} \`e una connettiva: \`e una famiglia di proposizioni indicizzata da una coppia di termini.

\section*{Cosa \emph{non} c'\`e: la dimensione costruttiva}

Tutto quello che abbiamo letto come logica \`e \textbf{intuizionista}. Significa, in concreto:

\begin{itemize}
\item \textbf{Terzo escluso non derivabile.} Il tipo $A \uplus \neg A$, in generale, \emph{non} \`e abitato da niente che si possa costruire uniformemente in $A$. Sarebbe una funzione $(A : \Type) \to A \uplus \neg A$ --- non c'\`e modo di scriverla in MLTT senza un postulato esterno. Logicamente: $A \vee \neg A$ non \`e un teorema dello schema F-I-E-C.
\item \textbf{Doppia negazione non eliminabile.} $\neg\neg A \to A$, in generale, non \`e abitato. Una prova di $\neg\neg A$ \`e una funzione $(A \to \bot) \to \bot$: ti dice che ``non si pu\`o produrre una funzione $A \to \bot$''. Da questa informazione \emph{non} si estrae un abitante di $A$. La differenza ``non c'\`e contro-esempio'' / ``c'\`e un esempio'' \`e costitutiva.
\end{itemize}

\noindent\textbf{Osservazione, non difetto.} La logica che affiora qui non \`e una logica \emph{mancante di qualcosa}: \`e una logica \emph{diversa}, dove la nozione di prova \`e pi\`u stringente. Una prova di $A$ \emph{esibisce} un abitante; una prova di $\exists x.\,P\,x$ \emph{esibisce} il testimone; una prova di $A \vee B$ \emph{decide} quale dei due vale. Per recuperare la matematica classica si pu\`o postulare LEM o DNE come ceri aggiuntivi (Stump lo fa esplicitamente in alcuni contesti); ma quello \`e un gesto separato, e va segnato come tale.

\section*{Tabella riassuntiva}

\begin{center}
\renewcommand{\arraystretch}{1.2}
\begin{tabular}{lll}
\toprule
\textit{Logica intuizionista} & \textit{Tipo} & \textit{Da quale F-I-E-C} \\
\midrule
$A \Rightarrow B$       & $A \to B$              & $\Pi$ non dip.\ (\S\ref{cap:pi}) \\
$A \wedge B$            & $A \times B$           & $\Sigma$ non dip.\ (\S\ref{cap:sigma}) \\
$\bot$ (falsit\`a)      & $\bot$                 & questa sezione, $\bot$ \\
$\neg A$                & $A \to \bot$           & definizione, $\to$ + $\bot$ \\
$\forall x : A.\,P\,x$  & $(x : A) \to P\,x$     & $\Pi$ dip.\ (\S\ref{cap:pi}) \\
$\exists x : A.\,P\,x$  & $\Sigma\,x \in A,\,P\,x$ & $\Sigma$ dip.\ (\S\ref{cap:sigma}) \\
$A \vee B$              & $A \uplus B$           & questa sezione, \S\ref{sec:coproduct} \\
$a = b$ in $A$          & $a \identita b$        & tipo identit\`a (\S\ref{cap:identita}) \\
\midrule
\multicolumn{3}{c}{\textit{non derivabili in generale}} \\
\midrule
LEM: $A \vee \neg A$    & $A \uplus (A \to \bot)$ & --- \\
DNE: $\neg\neg A \Rightarrow A$ & $((A \to \bot) \to \bot) \to A$ & --- \\
\bottomrule
\end{tabular}
\end{center}

\section*{Curry-Howard, come osservazione}

\noindent {\spacedlowsmallcaps{Quello che abbiamo appena descritto ha un nome}}. La coincidenza fra struttura dei tipi e struttura della logica intuizionista del primo ordine si chiama \textbf{corrispondenza di Curry-Howard} (talvolta \emph{Curry-Howard-de Bruijn}, o letta come ``isomorfismo proposizioni-come-tipi''). La presentano cos\`i:

\medskip
\begin{center}
proposizioni $\;\longleftrightarrow\;$ tipi \qquad prove $\;\longleftrightarrow\;$ termini
\end{center}
\medskip

\noindent Nei testi di logica dei tipi questa \`e quasi sempre l'apertura: ``decidiamo che tipi sono proposizioni, e dimostrare significa abitare''. Qui no. \`E arrivata alla fine, come constatazione di una corrispondenza gi\`a \emph{verificata caso per caso} nelle sezioni precedenti, e non come postulato di partenza.

\medskip

\noindent\textbf{Perch\'e la differenza importa.} Dichiarare la corrispondenza in apertura sarebbe un cero: ``assumiamo che tipi e proposizioni siano la stessa cosa''. Riconoscerla in chiusura \`e un'osservazione: ``avendo costruito i tipi per ragioni di calcolo (descrivere funzioni dipendenti, induzione, coppie), ci troviamo davanti che \emph{erano gi\`a} le regole della logica intuizionista''. Il nome resta utile come riferimento bibliografico; il \emph{contenuto} \`e tutto ci\`o che abbiamo gi\`a fatto.

\begin{oss}[Cero che non c'\`e]
La corrispondenza di Curry-Howard non figura nell'elenco dei ceri (\S\ref{cap:ceri}). Non perch\'e ce ne siamo dimenticati: perch\'e non \`e un'assunzione di MLTT. \`E una propriet\`a osservabile delle regole che MLTT ha gi\`a per altri motivi.
\end{oss}

% ---------------------------------------------------------------------
\chapter{Tipi notevoli}\label{cap:notevoli}

\epigraph{Come i limiti notevoli in analisi: poche costruzioni canoniche, da imparare una volta, da riconoscere ovunque.}{}

\noindent {\spacedlowsmallcaps{Un kit di attrezzi}}. Prima di entrare nella prova reale di Stump (\S\ref{cap:and-comm}), raccogliamo le costruzioni che si incontrano dappertutto nelle prove di Verified Functional Programming in Agda, scritte due volte: una alla Little Typer (con eliminatore esplicito), una in Agda (con pattern matching). \`E la palestra dove fissare l'\emph{esercizio centrale} di \S\ref{sec:stile}: leggere F-I-E-C sotto ogni clausola.

\bigskip

\noindent Cinque sezioni, in ordine di profondit\`a crescente: funzioni booleane (\S\ref{sec:not-funz-bool}), funzioni numeriche (\S\ref{sec:funz-nat}), propriet\`a del tipo identit\`a (\S\ref{sec:notevoli-id}), tipi piccoli di servizio (\S\ref{sec:tipi-piccoli}), operazioni su $\mathrm{Vec}$ (\S\ref{sec:notevoli-vec}). L'ultima sezione \`e il primo posto dove il cero~\ref{cero:pm} morde davvero.

\section{Funzioni booleane}\label{sec:not-funz-bool}

\subsection*{$\mathsf{not}$}

\paragraph{Definizione (Little Typer).} Eliminatore $\indB$ con motivo costante:
\[
\mathsf{not}\,b \;\definitorio\; \indB(b,\;\lambda\_.\,\mathbb{B},\;\ff,\;\tto)
\]

Letto: ``per costruire un $\mathbb{B}$ a partire da $b$, da' $\ff$ se $b = \tto$, $\tto$ se $b = \ff$''.

\paragraph{Definizione (Agda).}
\begin{agda}
not : 𝔹 → 𝔹
not tt = ff
not ff = tt
\end{agda}

Le due clausole \emph{sono} i due argomenti $c_t, c_f$ di $\indB$. Elaborazione triviale: niente matching su $\refl$, niente dot pattern.

\subsection*{$\_\&\&\_$ e $\_||\_$}

\paragraph{Definizione (Little Typer).}
\begin{align*}
b_1 \andd b_2 &\;\definitorio\; \indB(b_1,\;\lambda\_.\,\mathbb{B},\;b_2,\;\ff) \\
b_1 \,\|\, b_2 &\;\definitorio\; \indB(b_1,\;\lambda\_.\,\mathbb{B},\;\tto,\;b_2)
\end{align*}

Curry sul primo argomento: $b_2$ \`e nel contesto quando $\indB$ entra in scena, quindi sopravvive ai due rami.

\paragraph{Definizione (Agda).}
\begin{agda}
_&&_ : 𝔹 → 𝔹 → 𝔹
tt && b = b
ff && b = ff

_||_ : 𝔹 → 𝔹 → 𝔹
tt || b = tt
ff || b = b
\end{agda}

\begin{oss}[Cucitura con \S\ref{cap:and-comm}]
La definizione di \texttt{\_\&\&\_} appena vista \`e \textbf{esattamente} la stessa che useremo nel capitolo successivo per dimostrare \texttt{\&\&-comm}. Le ``regole di computazione di \texttt{\&\&}'' di cui parler\`a \S\ref{cap:and-comm} sono i due passi $\beta$ di $\indB$ applicato a $\tto$ e $\ff$.
\end{oss}

\subsection*{$\mathsf{if\_then\_else\_}$}

\paragraph{Definizione (Little Typer).} Per qualsiasi tipo $A$:
\[
\mathsf{if}\,b\,\mathsf{then}\,x\,\mathsf{else}\,y \;\definitorio\; \indB(b,\;\lambda\_.\,A,\;x,\;y)
\]

\paragraph{Definizione (Agda).}
\begin{agda}
if_then_else_ : {A : Set} → 𝔹 → A → A → A
if tt then x else y = x
if ff then x else y = y
\end{agda}

\noindent\textbf{Da fissare.} Le quattro funzioni booleane di base sono \emph{quattro istanze} dello stesso eliminatore $\indB$, distinte solo dal motivo e dai due rami. Nessuna nuova regola.

\section{Funzioni numeriche}\label{sec:funz-nat}

Stesso gesto su $\mathbb{N}$, ma ora i passi induttivi ricevono l'ipotesi induttiva.

\subsection*{$\_+\_$}

\paragraph{Definizione (Little Typer).} Induzione sul primo argomento, motivo $\mathbb{N} \to \mathbb{N}$ (cio\`e: ``$\mathsf{add}\,n$ \`e una funzione $\mathbb{N} \to \mathbb{N}$''):
\[
\mathsf{add}\,n \;\definitorio\; \indN(n,\;\lambda\_.\,\mathbb{N}\to\mathbb{N},\;\lambda m.\,m,\;\lambda k.\,\lambda \mathit{rec}.\,\lambda m.\,\suc\,(\mathit{rec}\,m))
\]

\noindent\textbf{Lettura dei pezzi:}
\begin{itemize}
\item motivo $\lambda\_.\mathbb{N}\to\mathbb{N}$: per ogni $n$, $\mathsf{add}\,n$ \`e una funzione $\mathbb{N}\to\mathbb{N}$;
\item caso $\zero$: $\lambda m.\,m$ --- aggiungere zero \`e l'identit\`a;
\item passo $\suc$: ricevo $k$, l'ipotesi induttiva $\mathit{rec} : \mathbb{N}\to\mathbb{N}$ (cio\`e $\mathsf{add}\,k$), e produco $\lambda m.\,\suc\,(\mathit{rec}\,m)$.
\end{itemize}

\paragraph{Definizione (Agda).}
\begin{agda}
_+_ : ℕ → ℕ → ℕ
zero  + m = m
suc k + m = suc (k + m)
\end{agda}

La ``chiamata ricorsiva'' \texttt{k + m} nel secondo ramo \`e \textbf{esattamente} il $\mathit{rec}\,m$ del passo induttivo di sopra. Il typechecker accetta perch\'e la ricorsione \`e su \texttt{k}, strutturalmente pi\`u piccolo di \texttt{suc k} (\S\ref{cap:nat}).

\subsection*{$\_*\_$}

\paragraph{Definizione (Little Typer).}
\[
\mathsf{mul}\,n \;\definitorio\; \indN(n,\;\lambda\_.\,\mathbb{N}\to\mathbb{N},\;\lambda m.\,\zero,\;\lambda k.\,\lambda \mathit{rec}.\,\lambda m.\,m + \mathit{rec}\,m)
\]

\paragraph{Definizione (Agda).}
\begin{agda}
_*_ : ℕ → ℕ → ℕ
zero  * m = zero
suc k * m = m + (k * m)
\end{agda}

\begin{oss}[Pattern: ``funzione binaria via induzione sul primo argomento'']
Sia \texttt{\_+\_} sia \texttt{\_*\_} usano lo stesso schema: $\indN$ sul primo argomento, motivo $\mathbb{N}\to\mathbb{N}$, caso base = funzione su $m$, passo = funzione su $m$ che usa la ricorsione. La maggior parte delle funzioni $\mathbb{N}\to\mathbb{N}\to X$ in Stump cap.~2 segue questo schema. Riconoscerlo basta a leggere mezzo capitolo del libro.
\end{oss}

\section{Propriet\`a del tipo identit\`a}\label{sec:notevoli-id}

\noindent {\spacedlowsmallcaps{Qui entra in scena il cero~\ref{cero:pm}}}. Tutte e quattro le costruzioni che seguono sono casi particolari di $\Jelim$ (\S\ref{cap:identita}), e tutte e quattro le loro versioni Agda fanno matching su $\refl$.

\subsection*{$\mathsf{sym}$ --- simmetria}

\paragraph{Definizione (Little Typer).} Dato $p : x \identita y$, vogliamo $y \identita x$. Si applica $\Jelim$ con motivo $C\,y\,p \definitorio y \identita x$. Caso base ($p = \refl$, $y = x$): $C\,x\,\refl = x \identita x$, abitato da $\refl$.
\[
\mathsf{sym}\,p \;\definitorio\; \Jelim(p,\;\lambda y\,p.\,y \identita x,\;\refl)
\]

\paragraph{Definizione (Agda).}
\begin{agda}
sym : {A : Set} {x y : A} → x ≡ y → y ≡ x
sym refl = refl
\end{agda}

\noindent\textbf{Cosa fa il matching di $\refl$.} Quando Agda vede la clausola \texttt{sym refl = ...}, istanzia $y := x$ nel resto del contesto. Quell'istanziazione \`e ci\`o che il cero~\ref{cero:pm} chiama in causa: in MLTT pura l'unico modo di ``sapere che $y = x$'' \`e specificare il motivo di $\Jelim$ e abitarlo nel caso $\refl$. Qui Agda lo fa da s\'e.

\subsection*{$\mathsf{trans}$ --- transitivit\`a}

\paragraph{Definizione (Little Typer).} Dati $p : x \identita y$, $q : y \identita z$, vogliamo $x \identita z$. Si applica $\Jelim$ su $p$ con motivo $C\,y\,p \definitorio y \identita z \to x \identita z$. Caso base: $C\,x\,\refl = x \identita z \to x \identita z$, abitato dall'identit\`a $\lambda q.\,q$.
\[
\mathsf{trans}\,p\,q \;\definitorio\; \Jelim(p,\;\lambda y\,p.\,y \identita z \to x \identita z,\;\lambda q.\,q)\,q
\]

\paragraph{Definizione (Agda).}
\begin{agda}
trans : {A : Set} {x y z : A} → x ≡ y → y ≡ z → x ≡ z
trans refl q = q
\end{agda}

Una sola clausola: il matching su $p = \refl$ istanzia $y := x$, e il tipo da abitare diventa $x \identita z$, che \`e gi\`a $q$.

\subsection*{$\mathsf{cong}$ --- congruenza}

\paragraph{Definizione (Little Typer).} Data $f : A \to B$ e $p : x \identita y$, vogliamo $f\,x \identita f\,y$. Motivo: $C\,y\,p \definitorio f\,x \identita f\,y$. Caso base: $C\,x\,\refl = f\,x \identita f\,x$, abitato da $\refl$.
\[
\mathsf{cong}\,f\,p \;\definitorio\; \Jelim(p,\;\lambda y\,p.\,f\,x \identita f\,y,\;\refl)
\]

\paragraph{Definizione (Agda).}
\begin{agda}
cong : {A B : Set} {x y : A} (f : A → B) → x ≡ y → f x ≡ f y
cong f refl = refl
\end{agda}

\subsection*{$\mathsf{subst}$ --- trasporto}

\paragraph{Definizione (Little Typer).} Data una propriet\`a $P : A \to \Type$, una prova $p : x \identita y$, e $\mathit{px} : P\,x$, vogliamo $P\,y$. Motivo: $C\,y\,p \definitorio P\,x \to P\,y$. Caso base: $C\,x\,\refl = P\,x \to P\,x$, abitato da $\lambda\mathit{px}.\,\mathit{px}$.
\[
\mathsf{subst}\,P\,p\,\mathit{px} \;\definitorio\; \Jelim(p,\;\lambda y\,p.\,P\,x \to P\,y,\;\lambda \mathit{px}.\,\mathit{px})\,\mathit{px}
\]

\paragraph{Definizione (Agda).}
\begin{agda}
subst : {A : Set} (P : A → Set) {x y : A} → x ≡ y → P x → P y
subst P refl px = px
\end{agda}

\begin{oss}[Quattro variazioni sullo stesso tema]
$\mathsf{sym}$, $\mathsf{trans}$, $\mathsf{cong}$, $\mathsf{subst}$ sono \textbf{quattro istanze di $\Jelim$}, distinte solo dal motivo. Il fatto che in Agda si scrivano tutte come una o due clausole con $\refl$ a sinistra \`e l'effetto del cero~\ref{cero:pm}: l'elaborazione del pattern matching trova il motivo da sola. Non magia, lavoro di unificazione.
\end{oss}

\section{Tipi piccoli di servizio}\label{sec:tipi-piccoli}

\subsection*{$\top$ --- il tipo unitario}

F-I-E-C completa (un costruttore, nessuna ricorsione):

\paragraph{F.} $\top : \Type$.

\paragraph{I.} $\star : \top$ (un solo costruttore senza premesse).

\paragraph{E.}
\[
\inferrule{u : \top \\ C : \top \to \Type \\ c : C\,\star}{\mathsf{ind}\text{-}\top(u, C, c) : C\,u}
\]

\paragraph{C.} $\mathsf{ind}\text{-}\top(\star, C, c) \definitorio c$.

\medskip

In Agda:

\begin{agda}
data ⊤ : Set where
  ★ : ⊤
\end{agda}

\noindent\textbf{Usi.} $\top$ \`e il tipo che ``non porta informazione'': abita la posizione di \emph{vero} nella logica, e serve come ``riempitivo'' quando un motivo deve restituire qualcosa di banale (lo vedremo immediato in \S\ref{sec:notevoli-vec} per $\mathsf{head}$).

\subsection*{$\neg$ --- negazione (richiamo)}

Gi\`a vista in \S\ref{cap:logica}: $\neg A \coloneqq A \to \bot$. Niente F-I-E-C nuova, \`e una definizione. Vale la pena ricordarla qui perch\'e compare in pratica sempre quando si lavora con $\mathsf{Dec}$, che chiameremo se serve:
\[
\mathsf{Dec}\,A \;\coloneqq\; A \uplus \neg A
\]

\noindent ``$A$ \`e decidibile'' = c'\`e o una prova di $A$ o una refutazione di $A$, e si sa quale.

\section{Operazioni su $\mathrm{Vec}\,A\,n$}\label{sec:notevoli-vec}

\noindent {\spacedlowsmallcaps{Dove il pattern matching comincia a non essere zucchero}}. Le funzioni che seguono sono ancora gestibili a mano con $\indV$ (\S\ref{cap:vec}), ma una di esse --- $\mathsf{head}$ --- richiede un trucco sul motivo che il pattern matching nasconde con grazia. \`E un buon posto per vedere il cero~\ref{cero:pm} al lavoro su tipi induttivi non banali.

\subsection*{$\mathsf{head}$ e $\mathsf{tail}$ --- e perch\'e il caso $\nil$ non c'\`e}

\paragraph{Definizione (Agda).}
\begin{agda}
head : {A : Set} {n : ℕ} → Vec A (suc n) → A
head (cons a as) = a

tail : {A : Set} {n : ℕ} → Vec A (suc n) → Vec A n
tail (cons a as) = as
\end{agda}

\noindent\textbf{La cosa che salta all'occhio.} Manca il caso $\nil$. Eppure Agda accetta come esaustivo. Perch\'e?

Perch\'e l'indice $\suc\,n$ nel tipo dell'argomento \emph{esclude} $\nil$: $\nil$ ha tipo $\mathrm{Vec}\,A\,\zero$, non $\mathrm{Vec}\,A\,(\suc\,n)$ per alcun $n$. L'unificatore di Agda riconosce questo e impone esaustivit\`a sui soli costruttori compatibili con l'indice. Questo \`e \textbf{esattamente} il tipo di lavoro che il cero~\ref{cero:pm} chiama in causa.

\paragraph{Versione Little Typer.} Volendo scrivere $\mathsf{head}$ con $\indV$ puro, l'unico modo \`e definire una funzione \emph{totale} su $\mathrm{Vec}\,A\,n$ qualunque $n$ sia, scegliendo un motivo che restituisca $\top$ nel caso $\zero$ e $A$ nel caso $\suc$:
\[
C : (n : \mathbb{N}) \to \mathrm{Vec}\,A\,n \to \Type
\qquad
C\,\zero\,\_ \definitorio \top
\qquad
C\,(\suc\,k)\,\_ \definitorio A
\]

Allora:
\[
\mathsf{head}'\,vs \;\definitorio\; \indV(vs,\;C,\;\star,\;\lambda k.\,\lambda a.\,\lambda \mathit{as}.\,\lambda\mathit{rec}.\,a)
\]

Per ottenere $\mathsf{head}$ ``vera'' --- quella che pretende un vettore di lunghezza $\suc\,n$ --- si specializza $\mathsf{head}'$ al caso in cui l'indice \`e $\suc\,n$, dove il motivo riduce ad $A$:
\[
\mathsf{head}\,(vs : \mathrm{Vec}\,A\,(\suc\,n)) \;\definitorio\; \mathsf{head}'\,vs : A
\]

\begin{oss}[Cosa abbiamo guadagnato col pattern matching]
La versione Agda \texttt{head (cons a as) = a} \`e \emph{una riga}. La versione Little Typer richiede di inventare un motivo che usa $\top$ come ``segnaposto'' per il caso che non interessa, definire una funzione totale, e poi specializzarla. Il contenuto matematico \`e lo stesso, ma la versione Agda \`e \emph{infinitamente} pi\`u leggibile. \`E il guadagno reale del cero~\ref{cero:pm}: il prezzo (assumere K, o appoggiarsi all'unificazione di \texttt{--without-K}) compra leggibilit\`a sostanziale.
\end{oss}

\subsection*{$\_\mathbin{+\!\!+}\_$ --- concatenazione}

\paragraph{Definizione (Agda).}
\begin{agda}
_++_ : {A : Set} {m n : ℕ} → Vec A m → Vec A n → Vec A (m + n)
nil         ++ ys = ys
(cons a as) ++ ys = cons a (as ++ ys)
\end{agda}

\paragraph{Definizione (Little Typer).} $\indV$ sul primo vettore, motivo che parametrizza sul vettore di destinazione:
\[
C : (m : \mathbb{N}) \to \mathrm{Vec}\,A\,m \to \Type
\qquad
C\,m\,\_ \definitorio \mathrm{Vec}\,A\,n \to \mathrm{Vec}\,A\,(m + n)
\]

\noindent\textbf{Cosa sta lavorando.} Nel caso $\nil$ ($m = \zero$): $C\,\zero\,\nil = \mathrm{Vec}\,A\,n \to \mathrm{Vec}\,A\,(\zero + n)$. Per la \emph{regola di computazione di $\_+\_$} (\S\ref{sec:funz-nat}), $\zero + n \definitorio n$, quindi il tipo \`e $\mathrm{Vec}\,A\,n \to \mathrm{Vec}\,A\,n$ --- abitato da $\lambda ys.\,ys$. Nel caso $\cons\,a\,\mathit{as}$ ($m = \suc\,k$): tipo $\mathrm{Vec}\,A\,n \to \mathrm{Vec}\,A\,(\suc\,k + n) = \mathrm{Vec}\,A\,(\suc\,(k+n))$, abitato da $\lambda ys.\,\cons\,a\,(\mathit{rec}\,ys)$.

\medskip

\noindent\textbf{Da fissare.} La regola di computazione di $\_+\_$ (definitoria, vista in \S\ref{sec:funz-nat}) lavora \emph{dentro il tipo}. Il typechecker la usa silenziosamente per verificare che $\mathrm{Vec}\,A\,(\zero + n)$ e $\mathrm{Vec}\,A\,n$ siano lo stesso tipo. \`E il Cero~\ref{cero:dipendenza} (un termine appare dentro un tipo) in azione operativa.

\subsection*{$\mathsf{map}$}

\paragraph{Definizione (Agda).}
\begin{agda}
map : {A B : Set} {n : ℕ} → (A → B) → Vec A n → Vec B n
map f nil         = nil
map f (cons a as) = cons (f a) (map f as)
\end{agda}

\paragraph{Definizione (Little Typer).} $\indV$ con motivo $C\,n\,\_ \definitorio \mathrm{Vec}\,B\,n$ (non dipende dal vettore, solo dalla lunghezza, perch\'e $\mathsf{map}$ preserva la lunghezza per costruzione):
\[
\mathsf{map}\,f\,vs \;\definitorio\; \indV(vs,\;\lambda n\,\_.\,\mathrm{Vec}\,B\,n,\;\nil,\;\lambda k.\,\lambda a.\,\lambda \mathit{as}.\,\lambda\mathit{rec}.\,\cons\,(f\,a)\,\mathit{rec})
\]

\begin{oss}[Il motivo costante sul vettore]
Nei tre notevoli di Vec di sopra, il motivo dipende dalla \emph{lunghezza} ma quasi mai dal \emph{vettore} stesso. Questo \`e tipico delle funzioni che ``lavorano in parallelo'' (mappa, zip, ecc.) rispetto a quelle che ``leggono'' il vettore (head, tail, lookup). Vederlo nei motivi rende esplicito cosa una funzione fa.
\end{oss}

\section{Cosa portare nel capitolo successivo}

In ordine di necessit\`a per \S\ref{cap:and-comm}:

\begin{itemize}
\item La definizione di \texttt{\_\&\&\_} come $\indB$ sul primo argomento con motivo costante (\S\ref{sec:not-funz-bool}). \`E ci\`o che render\`a meccaniche le ``riduzioni di \texttt{\&\&}'' nel quadro F-I-E-C di \texttt{\&\&-comm}.
\item L'introduzione di $\refl$ del tipo identit\`a, $\Jelim$ in standby, e il fatto che nel capitolo successivo \emph{non} useremo $\mathsf{sym}/\mathsf{trans}/\mathsf{cong}/\mathsf{subst}$ --- la prova si chiude solo con $\refl$, senza matching essenziale su prove di uguaglianza. Significa che il cero~\ref{cero:pm} \emph{non} \`e in scena nel caso \texttt{\&\&-comm}; arriver\`a nei capitoli successivi quando entrer\`a $\mathbb{N}$.
\end{itemize}

% ---------------------------------------------------------------------
\chapter{Anatomia di una prova: \texttt{\&\&-comm} in Iowa Agda Library}\label{cap:and-comm}

\noindent {\spacedlowsmallcaps{Mettiamo a frutto tutto il libro per leggere una prova reale}}. Ogni cosa che succede qui poggia su F-I-E-C gi\`a visti. \texttt{\_\&\&\_} la conosciamo dal kit dei notevoli (\S\ref{sec:not-funz-bool}); $\refl$ e il tipo identit\`a da \S\ref{cap:identita}; l'eliminatore $\indB$ da \S\ref{cap:bool}. E il \emph{patto di lettura} stabilito in \S\ref{cap:prima}: ogni clausola Agda verr\`a tradotta in F-I-E-C, segnalando se la traduzione \`e meccanica o se chiama in causa il cero~\ref{cero:pm}. \textbf{Anticipazione, per questo capitolo}: \`e meccanica. Il cero~\ref{cero:pm} non entra in scena --- entrer\`a quando faremo prove su $\mathbb{N}$ (cap.~2 di Stump).

\bigskip

L'esercizio del capitolo~1 di \emph{Verified Functional Programming in Agda} (Stump, libreria IAL):

\begin{agda}
module ch01.Booleans where
open import bool         -- definisce 𝔹, _&&_
open import bool-thms    -- teoremi precostruiti
open import eq           -- definisce _≡_ e refl (è il tipo identità di sezione 11)

&&-comm : ∀ (b1 b2 : 𝔹) → b1 && b2 ≡ b2 && b1
&&-comm tt tt = refl
&&-comm tt ff = refl
&&-comm ff tt = refl
&&-comm ff ff = refl
\end{agda}

Lo srotoliamo in tre tempi: che cos'\`e il \textbf{tipo}, che cos'\`e la \textbf{definizione di \texttt{\&\&}}, che cosa fa il \textbf{pattern matching}.

\section{Il tipo, strato per strato}

Il tipo \`e puro zucchero per due $\Pi$ annidati. Tolto lo zucchero:

\begin{agda}
-- strato 3 (file originale): forall + parentesi raggruppate
&&-comm : ∀ (b1 b2 : 𝔹) → b1 && b2 ≡ b2 && b1

-- strato 2: forall e' zucchero per Pi con tipo annotato (sez. forall)
&&-comm :   (b1 b2 : 𝔹) → b1 && b2 ≡ b2 && b1

-- strato 1: parentesi raggruppate e' zucchero per Pi annidati (sez. pi-azione)
&&-comm :   (b1 : 𝔹) → ((b2 : 𝔹) → (b1 && b2) ≡ (b2 && b1))

-- strato 0: MLTT puro, due applicazioni della regola di formazione di Pi (sez. pi)
--   con il tipo del risultato finale che e' un tipo identita' (sez. identita)
\end{agda}

Quindi \texttt{\&\&-comm} deve abitare \textbf{due $\Pi$ annidati}. Curry dipendente (\S\ref{cap:pi-azione}): nel secondo $\Pi$, \texttt{b1} \`e gi\`a nel contesto. Il tipo del risultato finale, \texttt{b1 \&\& b2 $\equiv$ b2 \&\& b1}, \`e il \textbf{tipo identit\`a} del linguaggio (\S\ref{cap:identita}) --- non l'uguaglianza definitoria.

\section{La definizione di \texttt{\&\&} come regole di computazione}

In IAL \texttt{\_\&\&\_} \`e definita:

\begin{agda}
_&&_ : 𝔹 → 𝔹 → 𝔹
tt && b = b
ff && b = ff
\end{agda}

\noindent\`E lo stesso \texttt{\_\&\&\_} di \S\ref{sec:not-funz-bool}: $\indB$ sul primo argomento, motivo costante $\lambda\_.\mathbb{B}$, rami $b_2$ e $\ff$. Le due clausole \emph{sono regole di computazione} sotto mentite spoglie. Cio\`e, il typechecker riduce automaticamente:
\begin{align*}
\tto \andd b &\definitorio b \\
\ff \andd b &\definitorio \ff
\end{align*}

\begin{oss}[Punto cruciale]
\texttt{\&\&} matcha \emph{solo sul primo argomento}. Se il primo \`e una variabile non istanziata, non riduce. Quindi:

\smallskip
\begin{itemize}
\item $\tto \andd b_2$ riduce a $b_2$. \checkmark
\item $\ff \andd b_2$ riduce a $\ff$. \checkmark
\item $b_1 \andd b_2$ con $b_1$ variabile aperta: \textbf{non riduce}. $\times$
\item $b_1 \andd \tto$ con $b_1$ variabile aperta: \textbf{non riduce}. $\times$
\end{itemize}
\end{oss}

\section{Il pattern matching come tre cose insieme}\label{sec:tre-fiec}

Le quattro clausole \texttt{\&\&-comm tt tt = refl} ecc.\ nascondono \textbf{tre F-I-E-C distinti} che lavorano in sequenza.

\paragraph{(a) Eliminazione di $\mathbb{B}$ (\S\ref{cap:bool}), due volte annidata.}
Il \texttt{tt}/\texttt{ff} sui due argomenti \texttt{b1, b2} \`e zucchero per $\indB$ applicato due volte. Agda impone esaustivit\`a: i quattro casi $(\tto,\tto), (\tto,\ff), (\ff,\tto), (\ff,\ff)$ non sono una scelta, sono i quattro rami obbligatori della doppia eliminazione.

\paragraph{(b) Riduzione computazionale di \texttt{\&\&} (\S\ref{cap:and-comm}.2).}
Una volta che \texttt{b1} e \texttt{b2} sono costruttori concreti, le clausole di \texttt{\&\&} riducono i due lati del tipo identit\`a a forme normali. Questo accade \emph{durante il type checking}: il tipo che si deve abitare in ciascun caso \`e quello \emph{dopo} la riduzione.

\medskip

Vediamolo caso per caso. In ciascun caso il tipo da abitare diventa:

\begin{center}
\renewcommand{\arraystretch}{1.2}
\begin{tabular}{lll}
\toprule
\textit{Caso} & \textit{Tipo originale} & \textit{Dopo riduzione di \texttt{\&\&}} \\
\midrule
$(\tto,\tto)$ & $\tto \andd \tto \identita \tto \andd \tto$ & $\tto \identita \tto$ \\
$(\tto,\ff)$  & $\tto \andd \ff \identita \ff \andd \tto$   & $\ff \identita \ff$ \\
$(\ff,\tto)$  & $\ff \andd \tto \identita \tto \andd \ff$   & $\ff \identita \ff$ \\
$(\ff,\ff)$   & $\ff \andd \ff \identita \ff \andd \ff$     & $\ff \identita \ff$ \\
\bottomrule
\end{tabular}
\end{center}

\smallskip

($\tto \andd \ff$ riduce a $\ff$ per la prima clausola di \texttt{\&\&}; $\ff \andd \tto$ riduce a $\ff$ per la seconda; ecc.)

\paragraph{(c) Introduzione del tipo identit\`a (\S\ref{cap:identita}) via $\refl$.}
In ogni caso, dopo la riduzione, i due lati sono \textbf{definitoriamente uguali}. Quindi $\refl$ li abita (Regola~\ref{reg:ponte}, il ponte). Le quattro \texttt{refl} nel codice \emph{non} sono lo stesso \texttt{refl}: sono quattro istanze diverse della regola di introduzione di \S\ref{cap:identita}, applicate in contesti diversi.

\section{Perch\'e \emph{non} basta matchare solo \texttt{b1}}

Sembrerebbe pulito scrivere solo due clausole:

\begin{agda}
&&-comm tt b2 = ???
&&-comm ff b2 = ???
\end{agda}

Vediamo il caso \texttt{tt}. Il tipo da abitare \`e $\tto \andd b_2 \identita b_2 \andd \tto$.

\smallskip
\begin{itemize}
\item $\tto \andd b_2$ riduce a $b_2$ (clausola 1 di \texttt{\&\&}). \checkmark
\item $b_2 \andd \tto$ \textbf{non riduce}, perch\'e $b_2$ \`e una variabile. $\times$
\end{itemize}
\smallskip

Quindi il tipo, dopo tutte le riduzioni possibili, resta $b_2 \identita b_2 \andd \tto$. Per $\refl$ servirebbero due lati definitoriamente uguali. Non lo sono. La propriet\`a ``qualcosa \texttt{\&\&} \texttt{tt} $\identita$ qualcosa'' \`e \textbf{proposizionalmente} vera (esiste una prova) ma \textbf{non definitoriamente} vera (non riduce). $\refl$ non la abita.

\begin{oss}[Conclusione strutturale]
Per ``sbloccare'' la riduzione di $b_2 \andd \tto$, serve case-split anche su $b_2$. \textbf{Quattro casi sono strutturalmente obbligatori}, non un capriccio della sintassi.
\end{oss}

\section{Il quadro completo}

Cosa abbiamo fatto, in una riga: il tipo $(b_1\,b_2 : \mathbb{B}) \to b_1 \andd b_2 \identita b_2 \andd b_1$ \`e due $\Pi$ annidati (\S\ref{cap:pi}, \S\ref{cap:pi-azione}); per abitarlo servono $\lambda b_1 \to \lambda b_2 \to \text{prova}$; per costruire la prova in modo che $\refl$ (\S\ref{cap:identita}) basti, i due lati del $\identita$ devono essere definitoriamente uguali (\S\ref{cap:giudizi}); per farli ridurre, servono i costruttori concreti, ottenuti via eliminazione di $\mathbb{B}$ (\S\ref{cap:bool}) in entrambe le variabili; in ognuno dei quattro casi le regole di computazione di \texttt{\&\&} (\S\ref{cap:and-comm}.2) riducono i due lati a $\tto \identita \tto$ o $\ff \identita \ff$, e $\refl$ chiude.

\bigskip

\noindent\textbf{Tre F-I-E-C --- $\mathbb{B}$ (E), \texttt{\&\&} (C), $\_\identita\_$ (I) --- che si intrecciano in quattro righe.} Niente magia, niente trucchi: solo le quattro regole del libro.

\bigskip

\begin{oss}[Cosa cambier\`a su $\mathbb{N}$]
Le prove del cap.~1 di Stump cadono tutte in questa famiglia: case-split puro + riduzione + $\refl$. Su $\mathbb{N}$ (cap.~2 di Stump) entreranno in gioco \emph{due} cose nuove. Primo: il \emph{passo induttivo} vero --- il caso $\suc\,k$ richieder\`a di chiamare l'ipotesi induttiva esplicitamente, cio\`e la chiamata ricorsiva nella definizione della funzione che dimostra. L'eliminatore $\indN$ entrer\`a in scena con la sua ricorsione, non solo come case-split. Secondo: il \emph{matching su prove di uguaglianza} (\texttt{sym refl = refl} ecc., \S\ref{sec:notevoli-id}) e il cero~\ref{cero:pm} entreranno in scena, perch\'e su $\mathbb{N}$ molte propriet\`a richiedono $\mathsf{cong}$, $\mathsf{subst}$, riscritture. Tutto quello che abbiamo costruito qui rester\`a, e si aggiungeranno questi due livelli.
\end{oss}
% =====================================================================
% APPENDICE — Indice (ceri e decisioni di design)
% =====================================================================
% Lo switch \appendix vive in main.tex prima dell'\input di questo file.

\chapter{Ceri e decisioni di design}\label{cap:ceri}

\noindent {\spacedlowsmallcaps{Questo \`e l'indice strutturale del libro}}. Ogni cosa che il libro assume senza dimostrarla --- ogni cosa che si \`e dovuta accendere come cero --- \`e qui. Separata da: ogni cosa che il libro \emph{decide} senza essere costretto a decidere cos\`i --- ogni scelta di design che ci sta come ci starebbe un'altra. Confondere le due cose \`e met\`a del lavoro mancato.

\section*{Ceri (assunzioni)}

\begin{center}
\renewcommand{\arraystretch}{1.4}
\begin{tabular}{p{0.5cm}p{5.2cm}p{3cm}p{4.8cm}}
\toprule
\textit{N.} & \textit{Assunzione} & \textit{Dove vive} & \textit{Perch\'e regge} \\
\midrule
1 & Esiste un metalinguaggio con i quattro giudizi & meta & \`e la lavagna su cui scrivi tutto il resto \\
\ref{cero:contesti} & Esistono i contesti e l'operazione $\vdash$ & meta & senza, niente ragionamento sotto ipotesi n\'e tipi dipendenti \\
\ref{cero:comp} & Le regole di computazione sono definitorie & meta & rispecchiano la struttura, non si provano dall'interno \\
\ref{cero:dipendenza} & Un termine pu\`o apparire dentro un tipo & scelta fondativa & \`e \emph{la} decisione che definisce MLTT \\
\midrule
(limite) & Coerenza del sistema & meta-meta & non dimostrabile dall'interno --- G\"odel \\
\bottomrule
\end{tabular}
\end{center}

\bigskip

\section*{Decisioni di design (scelte esplicite, non ceri)}

\begin{itemize}
\item \textbf{Eliminatore induttivo invece del solo case-split} per tipi ricorsivi --- $\indN$ vs $\mathsf{case\text{-}\mathbb{N}}$ (\S\ref{cap:bool} \emph{vs} \S\ref{cap:nat}). Potere dimostrativo: senza l'ipotesi induttiva nel passo, su $\mathbb{N}$ non si dimostra quasi niente. Si \emph{sarebbe potuto} scegliere solo case-split, e si avrebbe avuto un sistema diverso, pi\`u debole.

\item \textbf{Indice invece di parametro} per $n$ in $\mathrm{Vec}\,A\,n$ (\S\ref{cap:vec}). Permette a \texttt{cons} di cambiare la lunghezza producendo $\mathrm{Vec}\,A\,(\suc\,n)$ da $\mathrm{Vec}\,A\,n$. Se $n$ fosse parametro, sarebbe fisso per tutti i costruttori e \texttt{cons} non potrebbe esistere.

\item \textbf{Argomenti impliciti con \texttt{\{\}}} (\S\ref{cap:pi-azione}). Pura ergonomia, ma operativamente importante: \`e il riconoscimento sintattico che certi argomenti vivono solo nella passata 1 (type checking). L'argomento c'\`e sempre, sotto la sintassi.

\item \textbf{$\forall$ come zucchero per $\Pi$} (\S\ref{cap:forall}). Leggibilit\`a. Non \`e una regola nuova; \`e $\Pi$ con il permesso di omettere annotazioni di tipo quando Agda pu\`o inferirle.

\item \textbf{Parentesi raggruppate} \texttt{(a b c : A)} come zucchero per $\Pi$ annidati (\S\ref{cap:pi-azione}). Stessa famiglia di scelte: comodit\`a, sotto resta la regola di formazione unaria.
\end{itemize}

\bigskip

\noindent Ogni voce sopra \`e un punto in cui \emph{avresti potuto scegliere diversamente}. Saperlo \`e l'obiettivo: non confondere mai un'assunzione con un teorema, n\'e uno zucchero con una regola primitiva.

\bigskip
\bigskip

\section*{Dove vive ciascuna cosa --- la tavola finale}

\begin{center}
\renewcommand{\arraystretch}{1.25}
\begin{tabular}{p{6cm}p{3cm}p{4cm}}
\toprule
\textit{Cosa} & \textit{Dove vive} & \textit{Si dimostra?} \\
\midrule
$\tto, \ff, \zero, \suc\,n, \refl, \nil, \cons, \mathsf{inj}_1, \mathsf{inj}_2, \mathsf{fz}, \mathsf{fs}$ & linguaggio & s\`i, sono termini (introdotti dalle regole I) \\
$\mathbb{B}, \mathbb{N}, \Pi, \Sigma, A \uplus B, \_\identita\_, \mathrm{Fin}\,n, \mathrm{Vec}\,A\,n, \bot$ & linguaggio & s\`i, sono tipi (introdotti dalle regole F) \\
$\indB, \indN, \indS, \mathsf{ind}\text{-}\uplus, \Jelim$, applicazione di $\Pi$ & metalinguaggio & no, sono regole (E) \\
Regole di computazione (Cero~\ref{cero:comp}) & metalinguaggio & no, sono definitorie (C) \\
$\definitorio$ definitorio (giudizio) & metalinguaggio & no, \`e un'affermazione \\
$\identita$ proposizionale (tipo $\_\identita\_$) & linguaggio & s\`i, si abita con $\refl$ \\
Coerenza del sistema & meta-meta & non dall'interno (G\"odel) \\
\bottomrule
\end{tabular}
\end{center}

\bigskip

\noindent Pi\`u si scende verso le fondamenta, pi\`u il terreno diventa silenzioso. A un certo punto c'\`e sempre un cero alla madonna --- e va bene cos\`i, \emph{purch\'e sia esplicito}.

\bigskip

\noindent Lo svolgimento concreto del libro di Stump --- capitolo per capitolo, F-I-E-C alla mano --- \`e in Parte~IV (\S\ref{cap:stump-1} e seguenti).


\end{document}
